\section{Graph Theory}

\subsection{Core Concepts}

\begin{figure}[ht]
\centering
\begin{tikzpicture}[scale=1, transform shape]
    \node[ndout] (v1) at (5,5) {$v_1$};
    \node[ndint] (v2) at (1,4) {$v_2$};
    \node[ndint] (v3) at (11,5) {$v_3$};
    \node[ndout] (v4) at (3,3) {$v_4$};
    \node[ndout] (v5) at (7,3) {$v_5$};
    \node[ndout] (v6) at (11,3) {$v_6$};
    \node[ndout] (v7) at (5,1) {$v_7$};
    \node[ndout] (v8) at (7,-1) {$v_8$};
    \draw[arint] (v2) to (v4);
    \draw[arout] (v1) to (v4);
    \draw[arout] (v1) to (v5);
    \draw[arout] (v4) to (v7);
    \draw[arout] (v6) to (v8);
    \draw[arout] (v5) to (v7);
    \draw[arout] (v7) to (v8);
    \draw[arint] (v3) to (v6);
    \draw[arlat, bend left] (v1) to (v6);
\end{tikzpicture}
\caption{Conditional Acyclic Directed Mixed Graph (CADMG).}
\label{fig:cadmg}
\end{figure}


\begin{Def}[Conditional directed mixed graphs (CDMG)]
    \label{def-cdmg}
    A \emph{conditional directed mixed graph (CDMG)} $G$---per definition---consists of two (disjoint) sets of
    vertices (also called nodes):
    \begin{enumerate}[label=\roman*.)]
        \item $J$, whose elements are called input nodes,
        \item $V$, whose elements are called output nodes,
    \end{enumerate}
%    such that $J \cup V \neq \emptyset$;
    and two (disjoint) sets of edges:
    \begin{enumerate}[resume,label=\roman*.)]
        \item $E \ins (J \cup V) \times V$ the set of directed edges,
        \item $L \ins V \times V/((v_1,v_2) \sim (v_2,v_1)) $, the set of bi-directed edges, 
        \item[]    with: $(v_1,v_2) \in L \, \implies\, v_1\neq v_2 \land (v_2,v_1) \in L$.
    \end{enumerate}
\end{Def}

\begin{Not}
    \label{not-cdmg}
    Let $G=(J,V,E,L)$ be a CDMG.
    We will write:
    \begin{enumerate}
        %\item $G(V|\doit(J))$ to represent the graph $(J,V,E,L)$, where $E$ and $L$ are kept implicit in this notation.
        \item $v \in G$ to mean $v \in J \cup V$,
        \item $v_1 \tuh v_2 \in G$ to mean $(v_1,v_2) \in E$,
        \item $v_1 \hut v_2 \in G$ to mean $(v_2,v_1) \in E$,
        \item $v_1 \huh v_2 \in G$ to mean $(v_1,v_2) \in L$,
        \item $v_1 \suh v_2 \in G$ to mean that either $v_1 \tuh v_2 \in G$ or $v_1 \huh v_2 \in G$,
        \item $v_1 \hus v_2 \in G$ to mean that either $v_1 \hut v_2 \in G$ or $v_1 \huh v_2 \in G$,
        \item $v_1 \sus v_2 \in G$ to mean that either $v_1 \tuh v_2 \in G$ or $v_1 \hut v_2 \in G$ or $v_1 \huh v_2 \in G$.
    \end{enumerate}
    The star stands for a placeholder to mean: ``arrowhead or tail''.
\end{Not}

\begin{Def}
  Let $G=(J,V,E,L)$ be a CDMG.
  \begin{enumerate}
    \item If $v_1 \sus v_2 \in G$ then we call $v_1$ and $v_2$ \emph{adjacent in $G$}.
    \item Edges of the form $v_1 \hut v_2$ or $v_1 \huh v_2$ are called \emph{into $v_1$}.  \\
          Edges of the form $v_1 \tuh v_2$ or $v_1 \huh v_2$ are called \emph{into $v_2$}. 
    \item Edges of the form $v_1 \tuh v_2$ or $v_2 \hut v_1$ are called \emph{out of $v_1$}.
  \end{enumerate}
\end{Def}

\begin{Rem}
With the notations \ref{not-cdmg} the restrictions in definition \ref{def-cdmg} mean that the nodes $j \in J$ will not have any arrowheads pointing towards them: $ j \hus v \notin G$. Nodes $j \in J$ can only point towards nodes $v \in V$: edges $j \tuh v$ are allowed. Furthermore, no two nodes in $J$ are adjacent.
\end{Rem}

\begin{Def}[Walks]
     Let $G=(J,V,E,L)$ be a CDMG and $v,w \in G$.
\begin{enumerate}
    \item A \emph{walk} from $v$ to $w$ in $G$ is a finite alternating sequence of adjacent nodes and edges 
      \[v=v_0, a_0,  v_1, \dots v_{n-1}, a_{n-1}, v_n=w\] 
%      \[v=v_0 \sus v_1 \sus  \cdots \sus v_{n-1} \sus v_n=w\]
        in $G$ for some $n \ge 0$, i.e.\ such that for every $k=0,\dots,n-1$ 
        we have that $a_k = (v_k, v_{k+1}) \in E \cup L$ or $a_k = (v_{k+1}, v_k) \in E$, and with end nodes $v_0=v$ and $v_n=w$.
        An example walk from $v_0$ to $v_3$ could look like:
        \[v_0 \tuh v_1 \hut  v_2 \huh v_3, \qquad\text{ with } \qquad  v_0 \tuh v_1, v_2 \tuh v_1 \in E,\,v_2 \huh v_3 \in L. \]
        The same node may appear multiple times in a walk.
        Also the \emph{trivial walk} consisting of a single node $v_0 \in G$ is allowed (if $v=w$).
        The walk is called \emph{into $v_0$} if $a_0 = v_0 \hus v_1$, and \emph{out of $v_0$} if $a_0 = v_0 \tuh v_1$.
        Similarly, it is called \emph{into $v_n$} if $a_{n-1} = v_{n-1} \suh v_n$ and \emph{out of $v_n$} if $a_{n-1} = v_{n-1} \hut v_n$.
    \item A \emph{directed walk} from $v$ to $w$ in $G$ is of the form:
        \[v=v_0 \tuh v_1 \tuh  \cdots \tuh v_{n-1} \tuh v_n=w,\]
        for some $n \ge 0$,
        where all arrowheads point in the direction of $w$ and there are no arrowheads pointing back.
   % \item[] Directed walks exclude the trivial walk per definition.
    \item A \emph{bi-directed walk} from $v$ to $w$ in $G$ is of the form:
        \[v=v_0 \huh v_1 \huh  \cdots \huh v_{n-1} \huh v_n=w,\]
        for some $n \ge 0$, where all edges are bi-directed.
    \item A \emph{collider walk} from $v$ to $w$ in $G$ is of the form:
        \[v=v_0 \suh v_1 \huh  \cdots \huh v_{n-1} \hus v_n=w,\]
        for some $n \ge 0$, where all nodes in between $v$ and $w$ have two arrowheads pointing towards them (a.k.a.\ collider).
        Note that for $n=1$ this reads: $v\sus w \in G$.
    \item A walk is called \emph{path} if no node occurs more than once.
    \item A \emph{bifurcation} from $v$ to $w$ in $G$ is a walk of the form:
        \[v=v_0 \hut v_1 \hut  \cdots \hut v_{k-1} \hus v_k  \tuh \cdots \tuh v_{n-1} \tuh v_n=w,\]
   %             \[v=v_0 \hut v_1 \hut  \cdots \hut v_{k-1} \huh v_k  \tuh \cdots \tuh v_{n-1} \tuh v_n=w,\]
        such that $v \ne w$, the walk contains both endnodes exactly once, 
        every node has at most one arrowhead pointing towards it, and 
        both endnodes have exactly one arrowhead pointing towards them. 
          If the edge $v_{k-1} \hus v_k$ is directed ($v_{k-1} \hut v_k$) then we say that the bifurcation has \emph{source} $v_{k}$.
  \end{enumerate}
\end{Def}



\begin{Def}[Family relationships]
     Let $G=(J,V,E,L)$ be a CDMG, $v,w \in V$ and $A \ins J \cup V$ a subset of nodes.
     We then define:
     \begin{enumerate}
         \item The set of \emph{parents} of $v$ in $G$: \[\Pa^G(v):=\lC w \in G\,|\, w \tuh v \in G \rC.\]
         \item[] The set of \emph{parents} of $A$ in $G$: \[\Pa^G(A):= \bigcup_{v \in A} \Pa^G(v).\]
         %\item[] \PF{maybe better for sets (???) :The set of \emph{parents} of $A$ in $G$: \[\Pa^G(A):= \lp \bigcup_{v \in A} \Pa^G(v)\rp \sm A.\]
         %    but then  if $v \tuh v \in E$ then $v \notin \Pa^G(\lC v\rC ) \neq \Pa^G(v) \ni v$ ... }
         \item The set of \emph{children} of $v$ in $G$: \[\Ch^G(v):=\lC w \in G\,|\, v \tuh w \in G \rC.\]
         \item[] The set of \emph{children} of $A$ in $G$: \[\Ch^G(A):= \bigcup_{v \in A} \Ch^G(v).\]
         \item The set of \emph{siblings} of $v$ in $G$: \[\Sib^G(v):=\lC w \in G\,|\, v \huh w \in G \rC.\]
         %\item[] The set of \emph{children} of $A$ in $G$: \[\Ch^G(A):= \bigcup_{v \in A} \Ch^G(v).\]
         \item The set of \emph{ancestors} of $v$ in $G$: 
             \[\Anc^G(v):=\lC w \in G\,|\,\exists \text{ directed walk: } w \tuh \cdots \tuh v \in G \rC.\]
             Note: $v \in \Anc^G(v)$.
         \item[] The set of \emph{ancestors} of $A$ in $G$: \[\Anc^G(A):= \bigcup_{v \in A} \Anc^G(v).\]
             Note: $A \ins \Anc^G(A)$.
         \item The set of \emph{descendants} of $v$ in $G$: 
             \[\Desc^G(v):=\lC w \in G\,|\,\exists \text{ directed walk: } v \tuh \cdots \tuh w \in G \rC.\]
             Note: $v \in \Desc^G(v)$.
         \item[] The set of \emph{descendants} of $A$ in $G$: \[\Desc^G(A):= \bigcup_{v \in A} \Desc^G(v).\]
             Note: $A \ins \Desc^G(A)$.
         \item The set of \emph{non-descendants} of $A$ in $G$:
             \[ \NonDesc^G(A) := (J \cup V) \sm \Desc^G(A).\]
         \item The \emph{strongly connected component} of $v$ in $G$:
             \[ \Sc^G(v) := \Anc^G(v) \cap \Desc^G(v).\]
             Note: $v \in \Sc^G(v)$.
         \item[] The (union of) \emph{strongly connected components} of $A$ in $G$:
             \[ \Sc^G(A) := \bigcup_{v \in A} \Sc^G(v).\]
             Note: $A \ins \Sc^G(A)$.
         \item The \emph{district} of $v$ in $G$:
             \[\Dist^G(v) := \lC w \in G\,|\, \exists \text{ bi-directed walk: } 
             v \huh v_1 \huh \cdots \huh v_{n-1} \huh w \in G \rC.\]
             Note: $v\in \Dist^G(v)$.
         \item[] The \emph{district} of $A$ in $G$:
             \[\Dist^G(A) := \bigcup_{v \in A} \Dist^G(v).\]
             Note: $A \ins \Dist^G(A)$.
         %\item The d-\emph{Markov blanket} of $v$ in $G$:
         %    \[\MBl^G_d(v):=\lC w \in G\,|\,\exists \text{ collider walk: } 
         %    v \suh v_1 \huh \cdots \huh v_{n-1} \hus w \in G \rC \sm \{v\}.\]
         %    Note: $\lp\Pa^G(v) \cup \Ch^G(v) \cup \Pa^G(\Ch^G(v))\rp \sm \{v\} \ins \MBl^G_d(v)$.
         %\item[] The d-\emph{Markov blanket} of $A$ in $G$:
         %    \[\MBl^G_d(A) := \bigcup_{v \in A} \MBl^G_d(v) \sm A.\]
         %\item The $\sigma$-\emph{Markov blanket} of $A$ in $G$:
         %    \[\MBl^G_\sigma(A):= \MBl^{G^\acy}_d(A), \]
         %    where $G^\acy$ is an acyclification of $G$, a notion defined below in Definition~\ref{def:cdg_acyclification}.
     \end{enumerate}
\end{Def}


\begin{Def}[Acyclicity]
    \label{def-acylic}
    A  CDMG  $G=(J,V,E,L)$  is called \emph{acyclic} if there does not exist 
    any non-trivial directed walk from $v$ to itself in $G$ for any node $v \in G$.
\end{Def}

\begin{Def}
    A  Conditional Directed Mixed Graph (CDMG)  $G=(J,V,E,L)$  is called:
   \begin{enumerate}
       \item Conditional Acyclic Directed Mixed Graph (CADMG) if $G$ is acyclic.
       \item Directed Mixed Graph (DMG) if $J = \emptyset$.
       \item Acyclic Directed Mixed Graph (ADMG) if $G$ is acyclic and $J = \emptyset$.
       \item Conditional Directed Graph (CDG) if $L = \emptyset$.
       \item Directed Graph (DG) if $J = \emptyset$ and $L = \emptyset$.
       \item Conditional Directed Acyclic Graph (CDAG) if $G$ is acyclic and $L = \emptyset$.
       \item Directed Acyclic Graph (DAG) if $G$ is acyclic, $J=\emptyset$ and $L = \emptyset$.
   \end{enumerate}
\end{Def}

\begin{Def}[Topological order]
    Let $G=(J,V,E,L)$ be a CDMG.
    A \emph{topological order} of $G$ is a total order $<$ of $J \cup V$ such that for all $v,w \in G$:
    \[ v \in \Pa^G(w) \; \implies \; v < w.\]
    Equivalently, it can be described as an indexing of the nodes $J \cup V = \{v_1,\dots,v_K\}$ where parents always precede their children.
\end{Def}

\begin{Lem}
        A CDMG  $G=(J,V,E,L)$ is acyclic if and only if it has a topological order.
\end{Lem}

\begin{Def}[Predecessors]
   Let $G=(J,V,E,L)$ be a CDMG and $<$ a total order of $J \cup V$.
   The set of \emph{predecessors} of $v$ in $G$ are:
   \[\Pred^G_<(v) : = \lC w\in G\,|\, w < v \rC.\]
   We also put:
   \[\Pred^G_\le(v) : = \lC w\in G\,|\, w < v \rC \cup \{v\}.\]
\end{Def}


\subsection{Operations on Graphs}

%\subsection{Interventions on Graphs}
\subsubsection{Hard Interventions on Graphs}


\begin{figure}[ht]
\centering
\begin{tikzpicture}[scale=1, transform shape]
    \node[ndout] (v1) at (5,5) {$v_1$};
    \node[ndint] (v2) at (1,4) {$v_2$};
    \node[ndint] (v3) at (11,5) {$v_3$};
    \node[ndout] (v4) at (3,3) {$v_4$};
    \node[ndout] (v5) at (7,3) {$v_5$};
    \node[ndout] (v6) at (11,3) {$v_6$};
    \node[ndint] (v7) at (5,1) {$v_7$};
    %\node[ndout, below right=1.5cm and 1cm of v7] (v8) {$v_8$};
    \node[ndout] (v8) at (7,-1) {$v_8$};
    \draw[arint] (v2) to (v4);
    \draw[arout] (v1) to (v4);
    \draw[arout] (v1) to (v5);
    %\draw[arout] (v4) to (v7);
    \draw[arout] (v6) to (v8);
    %\draw[arout] (v5) to (v7);
    \draw[arint] (v7) to (v8);
    \draw[arint] (v3) to (v6);
    \draw[arlat, bend left] (v1) to (v6);
\end{tikzpicture}
\caption{The CADMG from Figure \ref{fig:cadmg} after hard intervention on node $v_7$.}
\end{figure}

\begin{Def}[Hard intervention on CDMGs]\label{def:G_hard_intervention}
    Let $G=(J,V,E,L)$ be a CDMG and $W \ins J \cup V$ a subset of nodes.
  The \emph{intervened CDMG} w.r.t.\ $W$ of $G$ is the CDMG: 
  \[ G_{\doit(W)}:=(J_{\doit(W)},V_{\doit(W)},E_{\doit(W)},L_{\doit(W)}),\] 
  %\[G(V\sm W|\doit(W \cup J)):=G_{\doit(W)}:=(J_{\doit(W)},V_{\doit(W)},E_{\doit(W)},L_{\doit(W)}),\] 
  where:
  \begin{enumerate}[label=\roman*.)]
      \item $J_{\doit(W)}:= J \cup W$,
      \item $V_{\doit(W)}:= V \sm W$,
      \item $E_{\doit(W)}:= E \sm \lC v \tuh w \,|\, v \in G, w  \in W  \rC$,
      \item $L_{\doit(W)}:= L \sm \lC v \huh w\,|\, v \in G, w \in W \rC$,
  \end{enumerate}
  where we turn all nodes from $W$ into input nodes and remove all edges into nodes from $W$.
\end{Def}

\begin{Rem}
    If $G$ is acyclic then also $G_{\doit(W)}$ is acyclic and a topological order for $G$ is also one for $G_{\doit(W)}$.
\end{Rem}

\begin{Lem}[Hard interventions commute]\label{hard-interventions-commute}
    Let $G:=(J,V,E,L)$ be a CDMG and $W_1, W_2 \ins J \cup V$ two  subsets of nodes from $G$. 
      Then we have:
      \[ \lp G_{\doit(W_1)} \rp_{\doit(W_2)} = \lp G_{\doit(W_2)} \rp_{\doit(W_1)} =  G_{\doit(W_1 \cup W_2)}. \]
      %This shows that our notion:
      %\[ G(V \sm (W_1 \cup W_2) | \doit(W_1 \cup W_2 \cup J))   \]
      %is unambiguous.
\end{Lem}

The following proposition expresses the existence of a bifurcation with a source in terms of ancestral relations in intervened graphs.
\begin{Prp}\label{prp:bifurcations_alternative}
  Let $G = \lt J, V, E, L \rt$ be a CMDG. For $v, w, c \in V \cup J$:
  there exists a bifurcation between $v$ and $w$ in $G$ with source $c$ if and only if $v \ne w$ and $c \in \Anc^{G_{\doit(w)}}(v) \sm \{v\}$ and $c \in \Anc^{G_{\doit(v)}}(w) \sm \{w\}$.
\end{Prp}
\begin{proof}
  A bifurcation between $v$ and $w$ with source $c$ is a walk in $G$ of the form $v \hut \cdots \hut c \tuh \cdots \tuh w$, where both $v$ and $w$ appear exactly once on the walk.
  This shows ``$\implies$''. For the other implication, note that $c \in \Anc^{G_{\doit(w)}}(v) \sm \{v\}$ implies that there is a non-trivial directed path from $c$ to $v$ that does not pass through $w$. 
  Similarly, $c \in \Anc^{G_{\doit(v)}}(w) \sm \{w\}$ implies that there is a non-trivial directed path from $c$ to $w$ that does not pass through $v$. 
  The concatenation of the two paths $v \hut \cdots \hut c \tuh \cdots \tuh w$ is then a bifurcation between $v$ and $w$ with source $c$.
\end{proof}

\subsubsection{Node Splitting Interventions on Graphs}

\begin{figure}[ht]
\centering
\begin{tikzpicture}[scale=1, transform shape]
\begin{scope}
    \node[ndout] (v1) at (5,5) {$v_1$};
    \node[ndint] (v2) at (1,4) {$v_2$};
    \node[ndint] (v3) at (11,5) {$v_3$};
    \node[ndout] (v4) at (3,3) {$v_4$};
%    \node[ndhst] (v5) at (7,3) {\nodepart{upper}{$v_5^o$} \nodepart{lower}{$v_5^i$}};
    \node[ndout] (v5o) at (7,3.5) {$v_5^o$};
    \node[ndint] (v5i) at (7,2.5) {$v_5^i$};
    \node[ndout] (v6) at (11,3) {$v_6$};
%    \node[ndhst] (v7) at (5,1) {\nodepart{upper}{$v_7^o$} \nodepart{lower}{$v_7^i$}};
    \node[ndout] (v7o) at (5,1.5) {$v_7^o$};
    \node[ndint] (v7i) at (5,0.5) {$v_7^i$};
    \node[ndout] (v8) at (7,-1) {$v_8$};
    \draw[arint] (v2) to (v4);
    \draw[arout] (v1) to (v4);
    \draw[arout] (v1) to (v5o);
    \draw[arout] (v4) to (v7o);
    \draw[arout] (v6) to (v8);
    \draw[arint] (v5i) to (v7o);
    \draw[arint] (v7i) to (v8);
    \draw[arint] (v3) to (v6);
    \draw[arlat, bend left] (v1) to (v6);
\end{scope}
\end{tikzpicture}
\caption{The CADMG from Figure \ref{fig:cadmg}  after a node-splitting hard intervention on $v_5$ and $v_7$.}
%Despite their different representation, the blue half nodes (here: $v_5^o,v_7^o$) are treated just as output nodes, and the half teal colored half nodes (here: $v_5^i,v_7^i$) are treated just as input nodes.}
\end{figure}



In this subsection we introduce \emph{node-splitting hard interventions}. 
They were introduced for the purpose of representing \emph{single-world intervention graphs (SWIGs)},
which represent the same output variable both before and after a hard intervention, see \cite{RR13a, RR13b}.
%The graphs were plotted using the library tikz-for-swigs \cite{Ric18swigtex}.

%\Joris{Restricted the definition to CDAMGS. This is because if one would like to equip simple SCMs describing a system at equilibrium with a semantics like this, the definition would need to be different.  It would need to work on the entire scc(s) of the intervened variable(s).}
\begin{Def}[Node-splitting hard intervention on CADMGs]
  \label{def:G_node-splitting_intervention}
    Let $G=(J,V,E,L)$ be a CADMG and $W \ins V$ a subset of the output nodes.
    The \emph{single-world intervention graph (SWIG)} w.r.t.\ $W$ of $G$ is the CADMG: 
    \[ G_{\swig(W)} :=\lp J_{\swig(W)}, V_{\swig(W)}, E_{\swig(W)},L_{\swig(W)} \rp,\]
    %\[ G_{\swig(W)} := G':=(J',V',E',L'),\]
    %\[G((V\sm W) \dcup W^o|\doit(W^i \dcup J)):=G':=(J',V',E',L'),\]
    constructed as follows.
    We first make two disjont copies of the nodes in $W$:
    \[ W^o:=\lC w^o\st w \in W \rC, \qquad W^i:=\lC w^i \st w \in W \rC.  \]
    Note that we consider $w^o \neq w^i$ for $w \in W$.
  However, for brevity, for $v \in J \cup V \sm W$ we put:
  \[ v^o:=v^i:=v.  \]
  We then define:
  \begin{enumerate}[label=\roman*.)]
      \item $J_{\swig(W)} := J \dcup W^i$,
      \item $V_{\swig(W)} := (V \sm W) \dcup W^o$, %which could be identified with $V$ again if we want to make the identification $W=W^o$.
      \item $E_{\swig(W)} := \lC v^i_1 \tuh v_2^o \st v_1 \tuh v_2 \in E  \rC$,
      \item $L_{\swig(W)} :=\lC v_1^o \huh v_2^o \st v_1 \huh v_2 \in L \rC$.
  \end{enumerate}
  where we turn all nodes of $W^i$ into input nodes, removing all edges into $W^i$,
  and we turn all nodes of $W^o$ into output nodes, removing all edges out of $W^o$.
\end{Def}

%\begin{Rem}
%    \Joris{Why would one even want to define this in the first place?}\PF{remove if irritating.}
%    If we wanted to also define node splitting for input nodes $j \in J$, then the only consistent way to do this is to
%replace $j$ by \emph{two input} nodes $j^o$ and $j^i$, where $j^i$ inherits all edges and $j^o$ has no edges.
%\end{Rem}

\begin{Rem}
    For a CADMG $G=(J,V,E,L)$, also $G_{\swig(W)}$ is acyclic.
    If $<$ is any topological order of $G$ given by enumerating all nodes $v \in J \cup V$ via:
    \[ v_1 < v_2 < \cdots < v_n,\]
    then, for instance, 
    a topological order for $G_{\swig(W)}$ can be achieved by assigning for a node $v_j \in W$ with index $j$ the node 
    $v_j^o$ the index $j-\frac{1}{3}$
    and $v_j^i$ the index $j+\frac{1}{3}$, and then ordering all nodes according to their index value.
\end{Rem}

\begin{Lem}[Two disjoint node-splitting hard interventions commute]
   Let $G=(J,V,E,L)$ be a CADMG and $W_1, W_2 \ins V$ two disjoint subsets of the output nodes from $G$. 
   Then the CADMG obtained from first node-splitting on $W_1$ and then node-splitting on $W_2$ is the same CADMG that arises from first node-splitting on $W_2$ and then node-splitting on $W_1$:
   \[ \lp G_{\swig(W_1)} \rp_{\swig(W_2)} =  \lp G_{\swig(W_2)} \rp_{\swig(W_1)}  =  G_{\swig(W_1 \dcup W_2)}.   \]
\end{Lem}


\begin{Lem}[Disjoint hard interventions and node-splitting hard interventions commute]
   Let $G=(J,V,E,L)$ be a CADMG and $W_1 \ins J \cup V$ and $W_2 \ins V$ two disjoint subsets of nodes from $G$. 
   Then the CADMG obtained from first hard intervening on $W_1$ and then node-splitting on $W_2$ is the same CADMG that arises from first node-splitting on $W_2$ and then hard intervening on $W_1$.
      \[ \lp G_{\doit(W_1)} \rp_{\swig(W_2)} =  \lp G_{\swig(W_2)} \rp_{\doit(W_1)}.   \]
\end{Lem}

\begin{Rem}
    Note that if $W_1$ and $W_2$ are not disjoint and $w \in W_1 \cap W_2 \ins V$ 
    then first hard intervening on $w$ turns $w$ into an input node, for now indicated as $w^i$, and a node-splitting hard intervention (if we would define it for input nodes) would not change $w^i$.
    If, on the other hand, we would first split the node $w$ into $w^o$ and $w^i$ then we would first need to resolve the ambiguity on which of those two the hard intervention should be applied. A hard intervention on $w^i$ would not do anything, but would leave the additional output node $w^o$ in the graph, while hard intervening on $w^o$ would turn $w^o$ into an additional input node, for now indicated as $(w^o)^i$. So in the latter case we are left with two input node $(w^o)^i$, which does not have any edges, and $w^i$, which might have outgoing edges.
\end{Rem}


\begin{comment}
\begin{figure}[ht]
    \centering
\begin{tikzpicture}[scale=1,transform shape]
\begin{scope}
    \node[ndout, xshift=0.0cm, yshift=0.0cm] (l) {$L$};
    \node[ndvst, below right=1.5cm and 0.4cm of l] (a) {
        					\nodepart{left}{$A$}
        					\nodepart{right}{$XXX$}};
    \node[ndout, right =1.5cm of a] (y) {$Y(a)$};
    \draw[arout] (l) to (a.north west);
    \draw[arout] (l) to (y);
    \draw[arint] (a) to (y);
\end{scope}
\end{tikzpicture}
    \caption{Conditional Acyclic Directed Mixed Graph (CADMG).}
\end{figure}
\end{comment}

\subsubsection{Intervention Nodes}

\begin{figure}[ht]
\centering
\begin{tikzpicture}[scale=1, transform shape]
    \node[ndout] (v1) at (5,5) {$v_1$};
    \node[ndint] (v2) at (1,4) {$v_2$};
    \node[ndint] (v3) at (11,5) {$v_3$};
    \node[ndout] (v4) at (3,3) {$v_4$};
    \node[ndout] (v5) at (7,3) {$v_5$};
    \node[ndout] (v6) at (11,3) {$v_6$};
    \node[ndout] (v7) at (5,1) {$v_7$};
    \node[ndout] (v8) at (7,-1) {$v_8$};
    \node[ndint] (Iv2) at (-1,4) {$I_{v_2}$};
    \node[ndint] (Iv7) at (3,1) {$I_{v_7}$};
    \draw[arint] (v2) to (v4);
    \draw[arout] (v1) to (v4);
    \draw[arout] (v1) to (v5);
    \draw[arout] (v4) to (v7);
    \draw[arout] (v6) to (v8);
    \draw[arout] (v5) to (v7);
    \draw[arout] (v7) to (v8);
    \draw[arint] (v3) to (v6);
    \draw[arint] (Iv7) to (v7);
    \draw[arlat, bend left] (v1) to (v6);
\end{tikzpicture}
  \caption{The CADMG from Figure \ref{fig:cadmg} with additional intervention node $I_{v_7}$ at node $v_7$. \JorisDoInput{Not clear why we would want to do this for input nodes. Suggest to not even define this operation for input nodes.}}
\end{figure}

More generally, interventions (both hard and soft) can be modeled graphically via auxiliary intervention nodes.

\begin{Def}[Extending CDMGs with intervention nodes]
  Let $G=(J,V,E,L)$ be a CDMG and $W \ins J \cup V$ a subset of nodes.
  The \emph{extended CDMG} of $G$ w.r.t.\ nodes $W \ins J \cup V$  and corresponding \emph{intervention nodes} 
  $I_W = \{I_w \,|\, w \in W\}$ %(with $I_j :=j$, for $j \in J$) 
  is the CDMG: 
  \[ G_{\doit(I_W)}:=(J_{\doit(I_W)},V_{\doit(I_W)},E_{\doit(I_W)},L_{\doit(I_W)}),\]
  \JorisDoInput{Suggest to change notation into $G_{\mathrm{int}(W)}$ instead.}
  %\[G(V|\doit(J \dcup I_W)):=G_{\doit(I_W)}=(J_{\doit(I_W)},V_{\doit(I_W)},E_{\doit(I_W)},L_{\doit(I_W)}),\]
  where:
  \begin{enumerate}[label=\roman*.)]
      \item $J_{\doit(I_W)}:= J\, \dot{\cup} \, \lC I_w\,|\, w \in W \rC$, %not W \sm J, makes weak and hard interventions commute
      \item $V_{\doit(I_W)}:= V$,
      \item $E_{\doit(I_W)}:= E\, \dot{\cup} \, \lC I_w \tuh w \,|\, w \in W\sm J \rC$,
      \item $L_{\doit(I_W)}:= L$,
  \end{enumerate}
  where we just add edges $I_w \tuh w$ for $ w \in W \sm J$, where $I_w$ represent intervention nodes.
  \JorisDoInput{Proposed alternative:
  \begin{enumerate}[label=\roman*.)]
    \item $J_{\doit(I_W)}:= (J \sm W)\, \dot{\cup} \, \lC I_w\,|\, w \in W \rC$, %not W \sm J, makes weak and hard interventions commute
    \item $V_{\doit(I_W)}:= V \cup (J \cap W)$,
      \item $E_{\doit(I_W)}:= E\, \dot{\cup} \, \lC I_w \tuh w \,|\, w \in W \rC$,
      \item $L_{\doit(I_W)}:= L$,
  \end{enumerate}
  where the former input nodes in $J \cap W$ become output nodes, and we add edges $I_w \tuh w$ for $ w \in W$ from the intervention nodes to their target nodes.
  }
%    \Joris{Why would you want to model a soft intervention on a subset of $J$? I don't see the point. Suggest to restrict to $W \subseteq V$ only!}\PF{have to check if we use $I_j$ anywhere for $j \in J$. see some discussion below.\\}
%  \PF{Alternative 0: We don't define $I_j$ for $j \in J$.\\}
%  \PF{Alternative 1: For $j \in J$ define $I_j\neq j$, but add NO edge $I_j \tuh j$ to the graph. This makes soft and hard interventions onto the same nodes commute.  (current version).\\}
%  \PF{Alternative 2: For $j \in J$ we make the identification $I_j:=j$. This is sometimes convenient for do-calculus rules.}
%  \Joris{I didn't find any do-calculus rule / part of its proof that makes use of this.}\\
%  \PF{Alternative 3: we could define $I_j$ for $j \in J$ with edge $I_j \tuh j$ and turn those $j$ into output nodes. This would give it some proper interpretataion as we would add a causal meachanism  for $j$ to the CBN.}
%  \Joris{Nice analysis. Perhaps there is no clear winner. I am fine with leaving it as-is because I don't have a convincing argument against it (even though personally would have gone for Alternative 0). Let's move on.}
\end{Def}

\begin{Rem}
    If a CDMG $G=(J,V,E,L)$ is acyclic then also $G_{\doit(I_W)}$ is acyclic and a topological order for
    $G_{\doit(I_W)}$ is also one for $G$.
    Any topological order of $G$ can be extended to one for $G_{\doit(I_W)}$, e.g.\ 
    by putting all the $I_w$ nodes first in the ordering.
\end{Rem}

\begin{Lem}[Adding intervention nodes commutes with disjoint hard interventions]
    Let $G=(J,V,E,L)$ be a CDMG and $W_1, W_2 \ins J \cup V$ two disjoint subsets of nodes from $G$.
      Then we have:
      \[ \lp G_{\doit(I_{W_1})} \rp_{\doit(I_{W_2})} = \lp G_{\doit(I_{W_2})} \rp_{\doit(I_{W_1})} 
      =  G_{\doit(I_{W_1 \cup W_2})}. \]
      We also have:
      \[ \lp G_{\doit(I_{W_1})} \rp_{\doit(W_2)} = \lp G_{\doit(W_2)} \rp_{\doit(I_{W_1})} =  G_{\doit(I_{W_1}, W_2)}. \]
   % So there is not ambiguity in the following notations and identifications:
   % \[ G(V|\doit(J\cup I_{W_1} \cup I_{W_2})) = G(V|\doit(J\cup I_{W_1 \cup W_2})), 
   % \qquad G(V\sm W_2|\doit(W_2 \cup J\cup I_{W_1})).\]
\end{Lem}

\begin{Lem}[Adding intervention nodes commutes with disjoint node-splitting hard interventions]
  Let $G=(J,V,E,L)$ be a CADMG and $W_1 \ins V$ and $W_2 \ins J \cup V$ two disjoint subsets of nodes from $G$.
    Then the CADMG that arises from first introducing intervention nodes $I_{W_2}$ and then splitting the nodes from $W_1$ is the same as the CADMG that arises from first splitting the nodes from $W_1$ and then introducing the intervention nodes $I_{W_2}$:
    \[  \lp G_{\swig(W_1)} \rp_{\doit(I_{W_2})} = \lp G_{\doit(I_{W_2})} \rp_{\swig(W_1)}. \]
  %  So there is no ambiguity in the notation:
  %  \[ G(V\sm W_1 \dcup W_1^o|\doit(W_1^i \dcup J \dcup I_{W_2})).  \]
\end{Lem}

\subsubsection{Marginalization of Graphs}

\begin{Def}[Marginalization a.k.a.\ latent projection on CDMGs]\label{def:G_marginalization}
  Let $G=(J,V,E,L)$ be a CDMG and $W \ins V$ a subset of output nodes.
  Then the \emph{marginalization of $G$ w.r.t.\ $W$} or the \emph{latent projection of $G$ onto $J\cup V \sm W$} is the CDMG:
  \[  G^{V\sm W | J} := G^{\sm W} := (J^{\sm W},V^{\sm W},E^{\sm W},L^{\sm W}), \] 
  %\[ G(V \sm W | \doit(J)) := G^{\sm W} = (J^{\sm W},V^{\sm W},E^{\sm W},L^{\sm W}), \] 
  where:
  \begin{enumerate}[label=\roman*.)]
      \item $J^{\sm W}:= J$,
      \item $V^{\sm W}:= V \sm W$,
      \item $E^{\sm W}$ consists of all directed edges $\ul{v} \tuh \ol{v}$ with $\ul{v},\ol{v} \in J \cup V \sm W$ 
          for which there exists a directed walk in $G$:
          \[\ul{v} \tuh w_1 \tuh \cdots \tuh w_{n-1} \tuh \ol{v}, \]
          where all intermediate nodes $w_1,\dots,w_{n-1} \in W$ (if any).\footnote{Note that this may introduce self-cycles.}
      \item $L^{\sm W}$ consists of all bi-directed edges $\ul{v} \huh \ol{v}$ with 
          $\ul{v},\ol{v} \in V \sm W$, $\ul{v} \neq \ol{v}$, 
          for which there exists a bifurcation in $G$:
          \[\ul{v} \hut w_1 \hut \cdots \hut w_{k-1} \hus w_k \tuh \cdots \tuh w_{n-1} \tuh \ol{v}, \]
          where all intermediate nodes $w_1,\dots,w_{n-1} \in W$ (if any).
      %\item $L^{\sm W}:= \lC \ul{v} \huh \ol{v}\,|\, \ul{v},\ol{v} \in V \sm W: 
      %    \exists \text{ bifurcation: } \ul{v} \hut w_1 \cdots w_{n-1} \tuh \ol{v} \in G 
      %    \text{ with } w_1,\dots,w_{n-1} \in W \rC$,
  \end{enumerate}
\end{Def}


\begin{Rem}\label{rem:marg_preserves_ancestors_bifurcations_acyclicity}
Marginalization preserves ancestral relations, bifurcations and acyclicity:
    \begin{enumerate}
        \item For $v_1,v_2 \in G$ with $v_1,v_2 \notin W$ we have the equivalence:
            \[ v_1 \in \Anc^G(v_2)\quad \iff \quad v_1 \in \Anc^{G^{\sm W}}(v_2).   \]
          \item For $v_1,v_2 \in G$ with $v_1,v_2 \notin W$ there is a bifurcation between $a$ and $b$ in $G$ if and only if there is a bifurcation between $a$ and $b$ in $G^{\sm W}$.\label{rem:marg_preserves_bifurcations}
        \item If the CDMG $G$ is acyclic then so is $G^{\sm W}$ and a topological order of $G$ induces a 
            topological order on $G^{\sm W}$ (by just ignoring the nodes from $W$).
    \end{enumerate}
\end{Rem}
\begin{proof}
  We prove \ref{rem:marg_preserves_bifurcations}.
  Let \[v_1=w_0 \hut w_1 \hut  \cdots \hut w_{k-1} \hus w_k  \tuh \cdots \tuh w_{n-1} \tuh w_n=v_2,\]
     be a bifurcation in $G$. If one marginalizes out a single node $u$ that is not on the bifurcation,
     then the same bifurcation exists in $G^{\setminus\{u\}}$. If one marginalizes out a single node
     $u$ that appears on the bifurcation (but that is not an endpoint) one obtains again a bifurcation in $G^{\setminus\{u\}}$. 
     The statement follows by induction.
\end{proof}

\begin{Lem}[Marginalizations commute]\label{marginalizations-commute}
      Let $G=(J,V,E,L)$ be a CDMG and $W_1, W_2 \ins V$ two disjoint subsets of output nodes. 
      Then we have:
      \[ \lp G^{\sm W_1} \rp^{\sm W_2} = \lp G^{\sm W_2} \rp^{\sm W_1} =  G^{\sm (W_1 \cup W_2)}. \]
     % This shows that the notation:
     % \[ G(V \sm (W_1 \cup W_2)|\doit(J))  \]
     % is unambiguous.
\end{Lem}

\begin{Lem}[Marginalization and intervention commute]
      Let $G=(J,V,E,L)$ be a CDMG and $W_1 \ins J \cup V$ and $W_2 \ins V$ two disjoint subsets of nodes from $G$.
      Then we have:
      \[ \lp G_{\doit(W_1)} \rp^{\sm W_2} = \lp G^{\sm W_2} \rp_{\doit(W_1)}. \]
     % This shows that the following notations are unambiguous:
     % \[G(V\sm (W_1 \cup W_2)|\doit(W_1 \cup J))=G_{\doit(W_1)}^{\sm W_2}.\]
      A similar statement holds for marginalizations and adding intervention nodes, and also for marginalizations and node-splitting interventions.
\end{Lem}

\begin{Lem}[Marginalizing out the output part of splitted nodes equals hard intervention]
     Let $G=(J,V,E,L)$ be a CDMG and $W \ins V$ be a subset of output nodes from $G$.
     Then the CDMG that arises by first splitting the nodes on $W$ 
      and then marginalizing out the nodes from $W^o$ 
     can be identified with the CDMG that arises by hard intervention on $W$:
     \[ G_{\doit(W)} \cong \lp G_{\swig(W)}\rp^{\sm W^o}, \qquad w \mapsto w^i. \]
\end{Lem}

\subsection{\textsigma-Separation}

\begin{Def}[Colliders and non-colliders]
  Let $G=(J,V,E,L)$ be a CDMG and $\pi$ a walk in $G$:
  \[ \pi =\lp  v_0 \sus \cdots \sus v_n \rp.  \]
  A node $v_k$, or more precisely, the position $k \in \lC 0,\dots, n \rC$,
  on the walk $\pi$ is called:
     \begin{enumerate}
     \item a \emph{non-collider} on $\pi$, if there is at most one arrowhead pointing towards $v_k$, i.e.\ if it falls into one of the following cases:
  \[\begin{array}{rccl}
    \text{ end-node: } & k \in \lC 0,n \rC, \\
    \text{ left chain: } & v_{k-1} \hut v_k \hus v_{k+1}, \\
    \text{ right chain: } & v_{k-1} \suh v_k \tuh v_{k+1},\\
    \text{ fork: } & v_{k-1} \hut v_k \tuh v_{k+1};
    \end{array}\]
  \item a \emph{collider} on $\pi$, if it is of the form: 
             \[v_{k-1} \suh v_k \hus v_{k+1},\] 
             i.e.\ if there are two arrowheads pointing towards $v_k$ on the walk $\pi$.
     \end{enumerate}
\end{Def}


\begin{Def}[Blockable and unblockable non-colliders]\label{def:unblockable_noncollider}
  Let $G=(J,V,E,L)$ be a CDMG and $\pi$ a walk in $G$:
  \[ \pi =\lp  v_0 \sus \cdots \sus v_n \rp.  \]
  We call a non-collider $v_k$ on $\pi$ an \emph{unblockable non-collider} on $\pi$ if it is not an end-node ($k \notin \{0,n\}$) and it only has outgoing edges on $\pi$ to nodes in the same strongly connected component of $G$. That is,
  it is one of the following patterns:
  \[\begin{array}{rccl}
    \text{ left chain: } & v_{k-1} \hut v_k \hus v_{k+1} & \text{ with } & v_{k-1} \in \Sc^G(v_{k})\\
    \text{ right chain: } & v_{k-1} \suh v_k \tuh v_{k+1} & \text{ with } & v_{k+1} \in \Sc^G(v_{k})\\
    \text{ fork: } & v_{k-1} \hut v_k \tuh v_{k+1} & \text{ with } & v_{k-1} \in \Sc^G(v_k) \land v_{k+1} \in \Sc^G(v_k)
  \end{array}\]
  Otherwise, $v_k$ is called a \emph{blockable non-collider} on $\pi$.
  This means that $v_k$ is either an end-node ($k \in \lC 0,n\rC$) or it has at least one outgoing arrow 
  $v_k \tuh v_{k\pm1}$ pointing to a node $v_{k\pm1}$ that lies in a different strongly connected component than $v_k$, i.e.\ $v_{k\pm1} \notin \Sc^G(v_k)$.
\end{Def}
  
\begin{Rem}
If $G$ is acyclic then all non-colliders are blockable.
\end{Rem}

\begin{comment}
\begin{Def}[$\sigma$-blocked walks]\label{def:sigma_blocking}
      Let $G=(J,V,E,L)$ be a CDMG and $C \ins J \cup V$ a subset of nodes and $\pi$ a walk in $G$:
      \[ \pi =\lp  v_0 \sus \cdots \sus v_n \rp.  \]
      \begin{enumerate}
          \item
      We say that the walk $\pi$ is \emph{$C$-$\sigma$-blocked} or \emph{$\sigma$-blocked by $C$} if:
      \begin{enumerate}[label=\roman*.)]
          \item $v_0 \in C$ or $v_n \in C$ or:
          \item there are two adjacent edges in $\pi$ of one of the following forms:
              \[\begin{array}{rccl}
  \text{ left chain: } & v_{k-1} \hut v_k \hus v_{k+1} & \text{ with } & v_k \in C \land v_k \notin \Sc^G(v_{k-1}),\\
  \text{ right chain: } & v_{k-1} \suh v_k \tuh v_{k+1} & \text{ with } & v_k \in C \land v_k \notin \Sc^G(v_{k+1}),\\
  \text{ fork: } & v_{k-1} \hut v_k \tuh v_{k+1} & \text{ with } & v_k \in C \land v_k \notin \Sc^G(v_{k-1}) \cap \Sc^G(v_{k+1}),\\
  \text{ collider: } & v_{k-1} \suh v_k \hus v_{k+1} & \text{ with } & v_k \notin \Anc^G(C).
              \end{array}\]
      \end{enumerate}
  \item We say that the walk $\pi$ is \emph{$C$-$\sigma$-open} if it is not $C$-$\sigma$-blocked.
  \end{enumerate}
\end{Def}
\end{comment}

\begin{comment}
      \Joris{Consider simplifying the above as:
        \begin{enumerate}[label=\roman*.)]
          \item $v_0 \in C$, or $v_n \in C$;
          \item or it contains a collider that is not in $\Anc^G(C)$:
              \[\begin{array}{rccl}
                \text{ collider: } & v_{k-1} \suh v_k \hus v_{k+1} & \text{ with } & v_k \notin \Anc^G(C);
              \end{array}\]
          \item or it contains a non-endpoint non-collider that is blockable and in $C$:
              \[\begin{array}{rccl}
  \text{ left chain: } & v_{k-1} \hut v_k \hus v_{k+1} & \text{ with } & v_k \in C \land v_k \notin \Sc^G(v_{k-1}),\\
  \text{ right chain: } & v_{k-1} \suh v_k \tuh v_{k+1} & \text{ with } & v_k \in C \land v_k \notin \Sc^G(v_{k+1}),\\
  \text{ fork: } & v_{k-1} \hut v_k \tuh v_{k+1} & \text{ with } & v_k \in C \land v_k \notin \Sc^G(v_{k-1}) \cap \Sc^G(v_{k+1}).\\
              \end{array}\]
      \end{enumerate}
          or even remove the equations altogether and express it in words (making use of the definitions of (blockable) (non)-collider), like we do in the remark below.
          }
  \PF{I don't see that this is simpler than the other versions we already have...}
\end{comment}

\begin{comment}
\begin{Rem}
    In words, the walk $\pi$ is:
  \begin{enumerate}
      \item   \emph{$C$-$\sigma$-open}  if and only if \emph{all}:
  \begin{enumerate}[label=\roman*.)]
      \item colliders $v_k$ on $\pi$ are in $\Anc^G(C)$, \emph{and}:
    \item blockable non-colliders $v_k$ on $\pi$ are \emph{not} in $C$.
      % all of its adjacent children, $v_{k\pm 1}$, on $\pi$ lie in its strongly connected component $\Sc^G(v_k)$. \Joris{the last part ``all of its neighboring children on $\pi$ lie in its strongly connected component'' is equivalent to ``unblockable'' (see Definition~\ref{def:unblockable_noncollider}); we might move that definition to here for convenience.}
  \end{enumerate}
      \item   \emph{$C$-$\sigma$-blocked} if and only if there \emph{exists}:
  \begin{enumerate}[label=\roman*.)]
    \item a collider $v_k$ on $\pi$ that is \emph{not} in $\Anc^G(C)$, \emph{or}:
    \item a blockable non-collider $v_k$ on $\pi$ in $C$.
  \end{enumerate}  \end{enumerate}
  Note that unblockable non-colliders are always $C$-$\sigma$-open, regardless of the subset $C \ins V \cup J$.
  Also note that if $G$ is acyclic then all non-colliders are blockable.
\end{Rem}
\end{comment}

\begin{Def}[$\sigma$-blocked walks]\label{def:sigma_blocking}
  Let $G=(J,V,E,L)$ be a CDMG and $C \ins J \cup V$ a subset of nodes and $\pi$ a walk in $G$:
  \[ \pi =\lp  v_0 \sus \cdots \sus v_n \rp.  \]
  We say that the walk $\pi$ is:
  \begin{enumerate}
    \item   \emph{$C$-$\sigma$-open} (or \emph{$\sigma$-open given $C$}) if and only if:
  \begin{enumerate}[label=\roman*.)]
      \item all colliders $v_k$ on $\pi$ are in $\Anc^G(C)$, \emph{and}:
    \item all blockable non-colliders $v_k$ on $\pi$ are \emph{not} in $C$.
      % all of its adjacent children, $v_{k\pm 1}$, on $\pi$ lie in its strongly connected component $\Sc^G(v_k)$. \Joris{the last part ``all of its neighboring children on $\pi$ lie in its strongly connected component'' is equivalent to ``unblockable'' (see Definition~\ref{def:unblockable_noncollider}); we might move that definition to here for convenience.}
  \end{enumerate}
      \item   \emph{$C$-$\sigma$-blocked}  (or \emph{$\sigma$-blocked given $C$}) if and only if:
  \begin{enumerate}[label=\roman*.)]
    \item there exists a collider $v_k$ on $\pi$ that is \emph{not} in $\Anc^G(C)$, \emph{or}:
    \item there exists a blockable non-collider $v_k$ on $\pi$ in $C$.
  \end{enumerate}  \end{enumerate}
  Note that unblockable non-colliders are always $C$-$\sigma$-open, regardless of the subset $C \ins V \cup J$.
\end{Def}


\begin{Def}[$\sigma$-separation]\label{def:sigma_separation}
    Let $G=(J,V,E,L)$ be a CDMG and $A,B,C \ins J \cup V$ (not necessarily disjoint) subset of nodes.
    We then say that:
    \begin{enumerate}
        \item \emph{$A$ is $\sigma$-separated from $B$ given $C$ in $G$}, in symbols:
    \[ A \Perp^\sigma_G B \given C, \]
    if every walk from a node in $A$ to a node in $J \cup B$ (sic!)\footnote{The choice to include $J$ here in this place is non-standard in the literature. However, if we include $J$ in this definition here the implied (asymmetric) separoid rules for d-/$\sigma$-separation will be of the same form as those for Markov kernels regarding conditional independence. This is the reason we include $J$ here.\label{fn:why-J}} is $C$-$\sigma$-blocked by $C$.
        \item If that property does not hold we will write:
            \[ A \nPerp^\sigma_G B \given C. \]
        \item We also define the special case:
                \[ A \Perp^\sigma_G B \qquad :\iff \qquad A \Perp^\sigma_GB \given \emptyset.\]
    \end{enumerate}
\end{Def}

The following result is often helpful to simplify proofs and making checking $\sigma$-separation feasible in practice.

\begin{restatable}{Prp}{restateprpsigmaopens}\label{prp:sigma_opens}
  Let $G=(J,V,E,L)$ be a CDMG. For $C \subseteq J \cup V$, and $w_1, w_2 \in J \cup V$, the following are
  equivalent:
  \begin{enumerate}
      \item there exists a $C$-$\sigma$-open \emph{path} between $w_1$ and $w_2$ in $G$;
      \item there exists a $C$-$\sigma$-open \emph{walk} between $w_1$ and $w_2$ in $G$;
      \item there exists a $C$-$\sigma$-open \emph{walk} between $w_1$ and $w_2$ in $G$ such that all its colliders lie in $C$ (and not just in $\Anc^G(C)$).
  \end{enumerate}
\end{restatable}

\begin{Rem}
    \begin{enumerate}
        \item By Proposition \ref{prp:sigma_opens} we have that $A \Perp^\sigma_G B \given C$ is equivalent to either of the following:
  \begin{enumerate}
      \item every walk from a node in $A$ to a node in $J \cup B$ is $C$-$\sigma$-blocked by $C$;
      \item every path from a node in $A$ to a node in $J \cup B$ is $C$-$\sigma$-blocked by $C$.
  \end{enumerate}
    \item Proposition \ref{prp:sigma_opens} also shows that if $A \nPerp^\sigma_G B \given C$ holds then:
  \begin{enumerate}
      \item there exists a (shortest) $C$-$\sigma$-open path from a node in $A$ to a node in $J \cup B$;
      \item there exists a (shortest) $C$-$\sigma$-open walk from a node in $A$ to a node in $J \cup B$ such that all its colliders lie in $C$.
  \end{enumerate}
  \end{enumerate}
  In practice we usually check if every path is $C$-$\sigma$-blocked or not. This is because there are, in contrast to walks, only a finite number of paths in a (finite) graph. In proofs, though, it often is easier to make use of walks, since these can be concatenated into walks (while one cannot in general concatenate two paths and again obtain a path).
\end{Rem}

\begin{Lem}[$\sigma$-separation under marginalization]\label{lem:stability_separation_marginalization}
        Let $G=(J,V,E,L)$ be a CDMG, $A,B,C \ins J \cup V$ and $D \ins V$ 
        be subsets of nodes such that:
        \[ D \cap \lp A \cup B \cup C\rp = \emptyset.\]
        Then we have the equivalence:
        \[ A \Perp^\sigma_G B \given C \qquad \iff \qquad  A \Perp^\sigma_{G^{\sm D} }B \given C. \]
\end{Lem}

\begin{Rem}
    If a CDMG $G$ is acyclic then all non-colliders are blockable.
    So, the partial condition for $\sigma$-separation
    ``a blockable non-collider in $C$''
    can be simplified to ``(any) non-collider in $C$''.
    
    So in the acyclic case we can simplify the notion of $\sigma$-separation, 
    which is usually referred to as \emph{d-separation}.
    However, in the non-acyclic setting d-separation (``(any) non-collider in $C$'') and $\sigma$-separation (``a blockable non-collider in $C$'') are clearly not equivalent anymore.

    It turned out that in the non-acyclic case $\sigma$-separation is the more general concept (and as said above it also captures the acyclic case equivalently well), see \cite{FM17,FM18,FM20,BFPM21}.
    We will first focus on CADMGs (acyclic) for which we can restrict ourselves to the somewhat simpler d-separation.
    Later, we will pick up $\sigma$-separation again when we deal with cycles.
\end{Rem}

\begin{comment}
\begin{Rem}
    If a CDMG $G$ is acyclic then all non-colliders are blockable,
     $\Sc^G(v)=\{v\}$ for all $v \in G$, and the additional conditions in ``$\sigma$-blocked'' 
    for the non-colliders are of the form $v_k \notin \Sc^G(v_{k\pm 1})=\{v_{k\pm1}\}$.
    These are then automatically satisfied, because the node $v_k$ is a parent of 
    the relevant other node $v_{k \pm1}$ and thus not equal to it.

    So in the acyclic case we can simplify the notion of $\sigma$-separation, which was then usually called \emph{d-separation}.
    However, in the non-acyclic setting d-separation and $\sigma$-separation are not equivalent anymore.
    It turns out that in the non-acyclic case $\sigma$-separation is the better concept (and as said above it also captures the acyclic case equivalently well), see \cite{FM17,FM18,FM20,BFPM21}.
    We will first focus on CADMGs (acyclic) for which we can restrict ourselves to the somewhat simpler d-separation.
    Later, we will pick up $\sigma$-separation again when we deal with cycles.
\end{Rem}
\end{comment}

\paragraph{Proofs - \textsigma-Open Walks and Paths}\hfill%

The following lemma will be convenient to relate $\sigma$-open walks and paths in the notion of $\sigma$-separation.


\begin{Lem}\label{lem:replace_walk}
Let $G=(J,V,E,L)$ be a CDMG, $C \subseteq V \cup J$ and $\pi = \lp  v_0 \sus \cdots \sus v_n \rp$ be a $C$-$\sigma$-open walk in $G$. 
  Suppose $v_i \in \Sc^G(v_j)$ for some $i,j \in \{0,\dots,n\}$ with $i < j$.
If we then replace the subwalk $v_i \sus \cdots \sus v_j$ of $\pi$ by 
  \begin{enumerate}[label=(\roman*)]
      \item a shortest directed path $v_i \tuh \cdots \tuh v_j$ in $G$ if $j = n$ or if $v_j \tuh v_{j+1}$ on $\pi$, or 
    \item a shortest directed path $v_i \hut \cdots \hut v_j$ in $G$ otherwise,
  \end{enumerate}
  then this new subwalk is entirely within $\Sc^G(v_j)$ and the modified walk $\pi'$ is still $C$-$\sigma$-open.
\begin{proof}
$\pi'$ cannot become $C$-$\sigma$-blocked at one of the initial nodes $v_0, \dots, v_{i-1}$ or at one of the final nodes $v_{j+1}, \dots, v_n$ on $\pi'$, since these nodes occur in the same local configuration on $\pi$ and are not $C$-$\sigma$-blocked on $\pi$ by assumption.
Furthermore, $\pi'$ cannot become $C$-$\sigma$-blocked at one of the nodes strictly between $v_i$ and $v_j$ on $\pi'$ (if there are any), since these nodes are all non-endnode non-colliders that only point to nodes in the same strongly connected component $\Sc^G(v_j)$. It is also worth noting that $\pi'$ cannot become $C$-$\sigma$-blocked at any of its endnodes, which could be $v_i$ or $v_j$ or both, because those are the same in $\pi$. So in the following we can w.l.o.g.\ assume that both $v_i$ and $v_j$ are non-endnodes of $\pi$ and thus $\pi'$.

Case (i). By assumption $v_j$ is either a fork or a right chain (or the right endnode) on $\pi$ that is $C$-$\sigma$-open. Since the same blocking criteria apply to $v_j$ on $\pi'$ it remains $C$-$\sigma$-open on $\pi'$. If $v_i=v_j$ then also $v_i$ is $C$-$\sigma$-open on $\pi'$ (if $v_i$ is the left endnode or not). If $v_i \neq v_j$, then the new directed path $v_i \tuh \cdots \tuh v_j$ in $\pi'$ is $C$-$\sigma$-open at $v_i$ because all nodes in between lie in the same stronly connected component $\Sc^G(v_i)$ (or $v_i$ is the left endnode anyways).

 Case (ii). Since case (i) is solved we can assume that we have $j<n$ with $v_j \hus v_{j+1}$ in $\pi$. 
 If $v_{i-1} \hut v_i$ on $\pi'$ (or $v_i$ the left endnode) then this case is analogous to case (i). 
 So we can also assume that we have $i > 0$ and $v_{i-1} \suh v_i$ on $\pi$. So $\pi$ looks as follows:
 \[\pi:\qquad \cdots v_{i-1} \suh v_i \sus \cdots \sus v_j \hus v_{j+1} \cdots. \]
 So there must be a smallest number $k \in \lC i,\dots,j\rC$ such that a collider appears at $v_k$ on $\pi$:
 \[\pi:\qquad \cdots v_{i-1} \suh v_i \tuh \cdots \tuh v_k \hus \cdots \sus v_j \hus v_{j+1} \cdots. \]
 Since $\pi$ is $C$-$\sigma$-open we have $v_k \in \Anc^G(C)$. Since $v_i \in \Anc^G(v_k)$ (otherwise $v_k$ would not be the first collider appearing after $v_i$) we thus have that also $v_i \in \Anc^G(C)$.
 So if we replace the subwalk $v_i \sus \cdots \sus v_j$ of $\pi$ by the shortest directed path $v_i \hut \cdots \hut v_j$ in $G$ we then get for $\pi'$ the following situation:
 \[\pi':\qquad \cdots v_{i-1} \suh v_i \hut \cdots \hut  v_j \hus v_{j+1} \cdots, \]
 which is then $C$-$\sigma$-open at $v_i$ as $v_i \in \Anc^G(C)$. Note that this holds also when $v_i=v_j$.
 If $v_i \neq v_j$ then $v_j$ is also $C$-$\sigma$-open on $\pi'$ as $v_j$ points left to a node in the same strongly connected component as $v_j$. 

 So in all cases $\pi'$ stays $C$-$\sigma$-open.
\end{proof}
\end{Lem}


\restateprpsigmaopens*
\begin{proof}
$3 \implies 2$ and $1 \implies 2$ are trivial. Note that paths are walks.

$2 \implies 3$: Suppose there exists a $C$-$\sigma$-open walk $\pi$ from $w_1$ to $w_2$.
Then consider a collider $v_{k-1} \suh v_k \hus v_{k+1}$ on $\pi$ with $v_k \in \Anc^G(C) \sm C$.
So there exists a non-trivial directed path from $v_k$ to a node $c_k \in C$ with all other nodes not in $C$.
If we then replace the collider at $v_k$ in $\pi$ by that path and its reverse we get:
\[ \cdots \sus v_{k-1} \suh v_k \tuh \cdots \tuh c_k \hut \cdots \hut v_k \hus v_{k+1} \sus \cdots.  \]
This walk is then $C$-$\sigma$-open at all places between $v_k$ on the left and $v_k$ on the right because they are non-colliders not in $C$. If we do this iteratively for all colliders not in $C$ we get the desired $C$-$\sigma$-open walk where all colliders lie in $C$.

$2 \implies 1$:
Let $\pi = (v_0 \sus \cdots \sus v_n)$ be a $C$-$\sigma$-open walk between nodes $v_0=w_1$ and $v_n=w_2$ in $G$. 
If a node $w$ occurs more than once on $\pi$, let $v_i$ be the first node on $\pi$ and $v_j$ be the last node on $\pi$ that are in $\Sc^G(w)$.
We now use Lemma~\ref{lem:replace_walk} to construct a new walk $\pi'$ from $\pi$ by replacing the subwalk between $v_i$ and $v_j$ of $\pi$ by a particular directed path in $\Sc^G(w)$ between $v_i$ and $v_j$ in such a way that $\pi'$ is still $C$-$\sigma$-open.
In $\pi'$, the number of nodes that occurs more than once is at least one less than in $\pi$, and all nodes within $\Sc^G(w)$ occur within a single segment.
This replacement procedure can be repeated until no nodes occur more than once. We have then obtained a
$C$-$\sigma$-open path between $w_1$ and $w_2$.
\end{proof}




\subsection{d-Separation}

\begin{Def}[d-blocked walks]\label{def:d-blocking}
      Let $G=(J,V,E,L)$ be a CDMG and $C \ins J \cup V$ a subset of nodes and $\pi$ a walk in $G$:
      \[ \pi =\lp  v_0 \sus \cdots \sus v_n \rp.  \]
      \begin{enumerate}
          \item
              We say that the walk $\pi$ is \emph{$C$-d-blocked} or \emph{d-blocked by $C$} to emphasize the use of the bi-directed edges.\footnote{The ``d'' here stands for ``directional''. d-separation was first only used for DAGs (without bi-directed edges).  For ADMGs it was then called m-separation in \cite{Ric03} But since the notion of m-separation is arguably the natural extension of d-separation and to avoid introducing more definitions, we will just call it d-separation as well, which will not create any ambiguity.} 
              if either:
      \begin{enumerate}[label=\roman*.)]
          \item $v_0 \in C$ or $v_n \in C$ or:
          \item there are two adjacent edges in $\pi$ of one of the following forms:
              \[\begin{array}{rccl}
           \text{ left chain: } & v_{k-1} \hut v_k \hus v_{k+1} & \text{ with } & v_k \in C,\\
           \text{ right chain: } & v_{k-1} \suh v_k \tuh v_{k+1} & \text{ with } & v_k \in C,\\
           \text{ fork: } & v_{k-1} \hut v_k \tuh v_{k+1} & \text{ with } & v_k \in C,\\
           \text{ collider: } & v_{k-1} \suh v_k \hus v_{k+1} & \text{ with } & v_k \notin \Anc^G(C).
              \end{array}\]
      \end{enumerate}
          \item We say that the walk $\pi$ is \emph{$C$-d-open} if it is not $C$-d-blocked.
      \end{enumerate}
\end{Def}

\begin{Rem}
          If we consider end-nodes, left chains, right chains and forks as \emph{non-colliders} then we can simply state:

      $\pi$ is d-blocked by $C$ if and only if it either contains a non-collider in $C$ or a collider not in $\Anc^G(C)$.
\end{Rem}

\begin{Def}[d-separation]
    Let $G=(J,V,E,L)$ be a CDMG and $A,B,C \ins J \cup V$ (not necessarily disjoint) subset of nodes.
    We then say that:
    \begin{enumerate}
        \item \emph{$A$ is d-separated from $B$ given $C$ in $G$}, in symbols:
    \[ A \Perp^d_G B \given C,  \]
    if every walk from a node in $A$ to a node in $J \cup B$ (sic!)\footref{fn:why-J} is $C$-d-blocked by $C$.
        \item If that property does not hold we will write:
            \[ A \nPerp^d_G B \given C.\]
    \end{enumerate}
\end{Def}

\begin{Rem}
    \begin{enumerate}
        \item A similar result from Proposition \ref{prp:sigma_opens} holds for d-separation as well.
        \item d-separation is stable under marginalization, similar to Lemma \ref{lem:stability_separation_marginalization}.
    \end{enumerate}
\end{Rem}


\begin{comment}
\subsection{\textsigma-Separation and d-Separation}

\PF{To get arround different notions of $\sigma$-separation we could only focus on the current $\sigma'$-separation version 
    ($\Anc^G(C)$ colliders). The important Lemma relating paths to walks could then be re-stated as the equivalence between the existence of:
\begin{enumerate}
    \item $C$-$\sigma'$-open path,
    \item $C$-$\sigma'$-open walk,
    \item $C$-$\sigma'$-open walk where all colliders are in $C$,  (<- only here comes the $C$-version in!!!).
\end{enumerate}
In proofs where shortes $C$-$\sigma$-open walks are required one could by this Lemma ask for:
``consider a shortest $C$-$\sigma'$-open walk where all colliders are in $C$.''\\
I think this would circumvent all problems with these two notions and could re-baptize $\sigma'$-separation as $\sigma$-separation.
}


\begin{Def}[$\sigma$-blocked walks]\label{def:sigma_blocking}
      Let $G=(J,V,E,L)$ be a CDMG and $C \ins J \cup V$ a subset of nodes and $\pi$ a walk in $G$:
      \[ \pi =\lp  v_0 \sus \cdots \sus v_n \rp.  \]
      \begin{enumerate}
          \item
      We say that the walk $\pi$ is \emph{$C$-$\sigma$-blocked} or \emph{$\sigma$-blocked by $C$} if either:
      \begin{enumerate}[label=\roman*.)]
          \item $v_0 \in C$ or $v_n \in C$ or:
          \item there are two adjacent edges in $\pi$ of one of the following forms:
              \[\begin{array}{rccl}
  \text{ left chain: } & v_{k-1} \hut v_k \hus v_{k+1} & \text{ with } & v_k \in C \land v_k \notin \Sc^G(v_{k-1}),\\
  \text{ right chain: } & v_{k-1} \suh v_k \tuh v_{k+1} & \text{ with } & v_k \in C \land v_k \notin \Sc^G(v_{k+1}),\\
  \text{ fork: } & v_{k-1} \hut v_k \tuh v_{k+1} & \text{ with } & v_k \in C \land v_k \notin \Sc^G(v_{k-1}) \cap \Sc^G(v_{k+1}),\\
  \text{ collider: } & v_{k-1} \suh v_k \hus v_{k+1} & \text{ with } & v_k \notin C.
              \end{array}\]
      \end{enumerate}
  \item  We say that the walk $\pi$ is \emph{$C$-$\sigma'$-blocked} or \emph{$\sigma'$-blocked by $C$} if the same rules as above hold, but where the collider case is replaced by:
              \[\begin{array}{rccl}
                  \text{ collider': } & v_{k-1} \suh v_k \hus v_{k+1} & \text{ with } & v_k \notin \Anc^G(C).
              \end{array}\]
  \item We say that the walk $\pi$ is \emph{$C$-$\sigma$-open} (\emph{$C$-$\sigma'$-open}, resp.) if it is not
      $C$-$\sigma$-blocked (not $C$-$\sigma'$-blocked, resp.).
  \end{enumerate}
\end{Def}

\begin{Rem}
    In words: the walk $\pi$ is \emph{$C$-$\sigma$-open} (\emph{$C$-$\sigma'$-open}, resp.)  if and only if:
  \begin{enumerate}
    \item both end-nodes of $\pi$ are not in $C$, and
    \item all colliders on $\pi$ are in $C$ (in $\Anc^G(C)$, resp.) and
    \item all (non-endpoint) non-colliders on $\pi$ are either not in $C$ or all of its neighboring children on $\pi$ lie in its strongly connected component. 
  \end{enumerate}
\end{Rem}


With the help of this lemma, we can prove that notions of $\sigma$-separation and $\sigma'$-separation are equivalent:

\begin{Prp}\label{prp:sigma_opens}
  Let $G=(J,V,E,L)$ be a CDMG. For $C \subseteq J \cup V$, and $v_1, v_2 \in J \cup V$, the following are
  equivalent:
  \begin{enumerate}
    \item there exists a $C$-$\sigma$-open walk between $v_1$ and $v_2$ in $G$;
    \item there exists a $C$-$\sigma'$-open walk between $v_1$ and $v_2$ in $G$;
    \item there exists a $C$-$\sigma'$-open path between $v_1$ and $v_2$ in $G$.
  \end{enumerate}
\end{Prp}



A simplified (and generally not eqivalent) version of the above is the following:

\begin{Def}[d-blocked walks]\label{def:d-blocking}
      Let $G=(J,V,E,L)$ be a CDMG and $C \ins J \cup V$ a subset of nodes and $\pi$ a walk in $G$:
      \[ \pi =\lp  v_0 \sus \cdots \sus v_n \rp.  \]
      \begin{enumerate}
          \item
              We say that the walk $\pi$ is \emph{$C$-d-blocked} or \emph{d-blocked by $C$}\footnote{The ``d'' here stands for ``directional''. d-separation was first only used for DAGs (without bi-directed edges).  For ADMGs it was then called m-separation in \cite{Ric03} to emphasize the use of the bi-directed edges. But since the notion of d-separation is the natural extension of d-separation, we will call it d-separation.} 
              if either:
      \begin{enumerate}[label=\roman*.)]
          \item $v_0 \in C$ or $v_n \in C$ or:
          \item there are two adjacent edges in $\pi$ of one of the following forms:
              \[\begin{array}{rccl}
           \text{ left chain: } & v_{k-1} \hut v_k \hus v_{k+1} & \text{ with } & v_k \in C,\\
           \text{ right chain: } & v_{k-1} \suh v_k \tuh v_{k+1} & \text{ with } & v_k \in C,\\
           \text{ fork: } & v_{k-1} \hut v_k \tuh v_{k+1} & \text{ with } & v_k \in C,\\
           \text{ collider: } & v_{k-1} \suh v_k \hus v_{k+1} & \text{ with } & v_k \notin C.
              \end{array}\]
      \end{enumerate}
      If we consider end nodes, left/right chains and forks as \emph{non-colliders} then we can simply say:\\
      $\pi$ is d-blocked by $C$ if it either contains a non-collider in $C$ or a collider not in $C$.
        \item  We say that the walk $\pi$ is \emph{$C$-d'-blocked} or \emph{d'-blocked by $C$} if the same rules as above hold, but where the collider case is replaced by:
              \[\begin{array}{rccl}
                  \text{ collider': } & v_{k-1} \suh v_k \hus v_{k+1} & \text{ with } & v_k \notin \Anc^G(C).
              \end{array}\]
          \item We say that the walk $\pi$ is \emph{$C$-d-open} (\emph{$C$-d'-open}, resp.) if it is not $C$-d-blocked (not $C$-d'-blocked, resp.), i.e.\ if all non-colliders are not in 
              $C$  and all colliders are in $C$ (in $\Anc^G(C)$, resp.).
      \end{enumerate}
\end{Def}
%\Joris{
%The following important property connects the two notions.
%\begin{Prp}
%Let $G=(J,V,E,L)$ be a CDMG, $C \ins J \cup V$ and $i, j \in V\cup J$. Then the following statements are equivalent:
%  \begin{enumerate}
%    \item there exists a $C$-$d$-open walk between $i$ and $j$ in $G$;
%    \item there exists a $C$-$d'$-open walk between $i$ and $j$ in $G$;
%    \item there exists a $C$-$d'$-open path between $i$ and $j$ in $G$.
%  \end{enumerate}
%\end{Prp}
%\begin{proof}
%TODO
%
%Since $C \subseteq \Anc^G(C)$, a collider not in $\Anc^G(C)$ is also not in $C$.
%Suppose there exists an $C$-$d$-open walk between $i$ and $j$.
%Then all non-colliders are not in $C$ and all colliders are in $C$. Hence all colliders are in $\Anc^G(C)$. Vice versa,
%suppose there exists an $C$-$d'$-open walk between $i$ and $j$. Then all non-colliders are not in $C$ and all
%colliders are in $\Anc^G(C)$. If a collider is not in $C$, we can extend the walk by replacing this collider 
%$\suh v \hus$ by the walk sequence
%$\suh v \tuh \dots \tuh c \hut \dots \hut v \hus$ with a directed path starting from $v$ and ending in the first
%node $c$ in $C$ it encounters, and from there tracing the same path backwards back to $v$. All nodes on this
%walk except for $c$ ($v$ included) are non-colliders and are not in $C$. Hence this walk is $C$-$d$-open.
%\end{proof}
%}

\begin{Def}[d-separation/$\sigma$-separation]
    Let $G=(J,V,E,L)$ be a CDMG and $A,B,C \ins J \cup V$ (not necessarily disjoint) subset of nodes.
    We then say that:
    \begin{enumerate}
        \item \emph{$A$ is $\sigma$-separated from $B$ given $C$ in $G$} (\emph{d-separated}, resp.), in symbols:
    \[ A \Perp^\sigma_G B \given C, \qquad A \Perp^d_G B \given C, \text{ resp.,}  \]
    if every walk from a node in $A$ to a node in $J \cup B$ (sic!) is $C$-$\sigma$-blocked by $C$ ($C$-d-blocked, resp.).
        \item \emph{$A$ is $\sigma'$-separated from $B$ given $C$ in $G$} (\emph{d'-separated}, resp.), in symbols:
            \[ A \Perp^{\sigma'}_G B \given C, \qquad A \Perp^{d'}_G B \given C, \text{ resp.,}  \]
            if every \emph{path} (sic!) from a node in $A$ to a node in $J \cup B$ (sic!) is $C$-$\sigma'$-blocked by $C$ ($C$-d'-blocked, resp.).
        \item If the corresponding property does not hold we will write:
            \[ A \nPerp^\sigma_G B \given C, \qquad A \nPerp^d_G B \given C, \qquad A \nPerp^{\sigma'}_G B \given C, \qquad A \nPerp^{d'}_G B \given C, \text{ resp.} \]
        \item For $\Perp=\Perp^\sigma_G$, $\Perp^d_G$, $\Perp^{\sigma'}_G$, $\Perp^{d'}_G$, resp., we also define the special case:
                \[ A \Perp B \qquad :\iff \qquad A \Perp B | \emptyset.\]
    \end{enumerate}
\end{Def}

\begin{Rem}
    If a CDMG $G$ is acyclic then $\Sc^G(v)=\{v\}$ for all $v \in G$ and the additional conditions in ``$\sigma$-blocked'' 
    for the non-colliders are of the form $v_k \notin \Sc^G(v_{k\pm 1})=\{v_{k\pm1}\}$.
    These are then automatically satisfied, because the node $v_k$ is a parent of 
    the relevant other node $v_{k \pm1}$ and thus not equal to it. 
    So in the acyclic case d-separation and $\sigma$-separation are equivalent.
    It turns out that in the non-acyclic case $\sigma$-separation is the better concept (and as said above it also captures the acyclic case equivalently well), see \cite{FM17,FM18,FM20,BFPM21}.
    We will first focus on CADMGs (acyclic) for which we can restrict ourselves to the somewhat simpler d-separation.
    Later, we will pick up $\sigma$-separation again when we deal with cycles.
\end{Rem}

\begin{Lem}[d-separation/$\sigma$-separation under marginalization]\label{lem:stability_separation_marginalization}
        Let $G=(J,V,E,L)$ be a CDMG, $A,B,C \ins J \cup V$ and $D \ins V$ 
        be subsets of nodes such that:
        \[ D \cap \lp A \cup B \cup C\rp = \emptyset.\]
        Then we have the equivalence:
        \[ A \Perp^d_G B \given C \qquad \iff \qquad  A \Perp^d_{G^{\sm D} }B \given C. \]
        We also have:
        \[ A \Perp^\sigma_G B \given C \qquad \iff \qquad  A \Perp^\sigma_{G^{\sm D} }B \given C. \]
\end{Lem}

%\begin{Def}[Acyclification of CDMGs]
%    \label{def-acyclification}
%Let $G=(J,V,E,L)$ be a CDMG.
%We define the \emph{acyclification} $G^\acy = (J^\acy,V^\acy,E^\acy,L^\acy)$ of $G$ as follows:
%\begin{enumerate}
%    \item  $ J^\acy:=J$ and $V^\acy :=V$.
%    \item  $v \tuh w \in E^\acy$ iff $v \notin \Sc^G(w)$ and there exists $\tilde w \in \Sc^G(w)$ such that $v \tuh \tilde w \in E$:
%        \[ E^\acy:=\lC v \tuh w \,|\, v \notin \Sc^G(w) \, \land\, \exists \tilde w \in \Sc^G(w):\; v \tuh \tilde w \in E \rC.  \]
%    \item  $v \huh w \in L^\acy$ iff $v \notin \Sc^G(w)$  and there 
%        exists $\tilde v \in \Sc^G(v)$ and $\tilde w \in \Sc^G(w)$ 
%        such that $\tilde v \huh \tilde w \in L$:
%        \[ L^\acy:=\lC v \huh w \,|\,v \notin \Sc^G(w) \, \land\, \exists \tilde v \in \Sc^G(v) \,\land\, \exists \tilde w \in \Sc^G(w):\; \tilde v \huh \tilde w \in L \rC.  \]
%\end{enumerate}
%In other words, the acyclification draws edges $ v \tuh w \in \Sc^G(w)$ from outside a strongly connected component to every node
%$\tilde w \in \Sc^G(w)$ of that strongly connected component, similarly for bi-directed edges, 
%and then removes all edges inside each strongly connected component of $G$. 
%Since strongly connected components can themselves ordered in an acyclic fashion the resulting graph $G^\acy$ is acyclic, thus a CADMG.
%\end{Def}
%
%\begin{Thm}
%    \label{thm-acy-sigma-sep}
%    Let $G=(J,V,E,L)$ be a CDMG and $G^\acy$ its acyclification.
%    For $A,B,C \ins J \cup V$ we then have the equivalence:
%    \[ A \Perp^\sigma_G B \given C  \qquad\iff\qquad A \Perp^\sigma_{G^\acy} B \given C \qquad\iff\qquad A \Perp^d_{G^\acy} B \given C.  \]
%\end{Thm}

\subsubsection{Alternative definition of $\sigma$-blocking}

Revisit Definition~\ref{def:sigma_blocking} of $\sigma$-blocking. In the literature, one also encounters
a slightly different notion with the same name. In this optional section, we will clarify the relationship
between the two. %For the rest of the course, we will keep using Definition~\ref{def:sigma_blocking}.
\begin{Def}\label{def:sigma_blocking_2}
  Let $G=\langle J,V,E,L \rangle$ be a CDMG, $C \ins J \cup V$ a subset of nodes and $\pi$ a walk (or path) in $G$:
  \[ \pi =\lp  v_0 \sus \cdots \sus v_n \rp.  \]
\begin{enumerate}
  \item
    We say that the walk (or path) $\pi$ is \emph{$C$-$\sigma'$-blocked} or \emph{$\sigma'$-blocked by $C$} if either:
      \begin{enumerate}[label=\roman*.)]
          \item $v_0 \in C$ or $v_n \in C$ or:
          \item there are two adjacent edges in $\pi$ of one of the following forms:
              \[\begin{array}{rccl}
  \text{ left chain: } & v_{k-1} \hut v_k \hus v_{k+1} & \text{ with } & v_k \in C \land v_k \notin \Sc^G(v_{k-1}),\\
  \text{ right chain: } & v_{k-1} \suh v_k \tuh v_{k+1} & \text{ with } & v_k \in C \land v_k \notin \Sc^G(v_{k+1}),\\
  \text{ fork: } & v_{k-1} \hut v_k \tuh v_{k+1} & \text{ with } & v_k \in C \land v_k \notin \Sc^G(v_{k-1}) \cap \Sc^G(v_{k+1}),\\
                \text{ collider: } & v_{k-1} \suh v_k \hus v_{k+1} & \text{ with } & v_k \notin \Anc^G(C).
              \end{array}\]
      \end{enumerate}
  \item We say that the walk (or path) $\pi$ is \emph{$C$-$\sigma'$-open} if it is not $C$-$\sigma'$-blocked. 
\end{enumerate}
\end{Def}
The only difference is in the condition for a collider node $v_k$: to $\sigma$-block
a walk one needs $v_k \not\in C$, whereas to $\sigma'$-block a walk it suffices that $v_k \not\in \Anc^G(C)$.

\end{comment}

\begin{comment}
\paragraph{Proofs - \textsigma-Open Walks and Paths}\hfill%

The following lemma will be convenient to relate the different notions of $\sigma$-separation.


\begin{Lem}\label{lem:replace_walk}
Let $G=(J,V,E,L)$ be a CDMG, $C \subseteq V \cup J$ and $\pi = \lp  v_0 \sus \cdots \sus v_n \rp$ be a $C$-$\sigma'$-open walk in $G$. 
  Suppose $v_i \in \Sc^G(v_j)$ for some $i,j \in \{0,\dots,n\}$ with $i < j$.
We can then replace the subwalk $v_i \sus \cdots \sus v_j$ of $\pi$ by 
  \begin{enumerate}[label=(\roman*)]
      \item a shortest directed path $v_i \tuh \cdots \tuh v_j$ in $G$ if $j = n$ or if $v_j \tuh v_{j+1}$ on $\pi$, or 
          %\PF{"if $\tuh \underbrace{v_j \sus v_{j+1}}_{\text{on }\pi}$ would be $C$-$\sigma'$-open"}
    \item a shortest directed path $v_i \hut \cdots \hut v_j$ in $G$ otherwise,
  \end{enumerate}
that is entirely within $\Sc^G(v_j)$ and such that the modified walk $\pi'$ is still $C$-$\sigma'$-open.
\begin{proof}
$\pi'$ cannot become $C$-$\sigma'$-blocked at one of the initial nodes $v_0, \dots, v_{i-1}$ or at one of the final nodes $v_{j+1}, \dots, v_n$ on $\pi'$, since these nodes occur in the same local configuration on $\pi$ and are not $C$-$\sigma'$-blocked on $\pi$ by assumption.
Furthermore, $\pi'$ cannot become $C$-$\sigma'$-blocked at one of the nodes strictly between $v_i$ and $v_j$ on $\pi'$ (if there are any), since these nodes are all non-endnode non-colliders that only point to nodes in the same strongly connected component $\Sc^G(v_j)$. It is also worth noting that $\pi'$ cannot become $C$-$\sigma'$-blocked at any of its endnodes, which could be $v_i$ or $v_j$ or both, because those are the same in $\pi$. So in the following we can w.l.o.g.\ assume that both $v_i$ and $v_j$ are non-endnodes of $\pi$ and thus $\pi'$.

Case (i). By assumption $v_j$ is either a fork or a right chain (or the right endnode) on $\pi$ that is $C$-$\sigma'$-open. Since the same blocking criteria apply to $v_j$ on $\pi'$ it remains $C$-$\sigma'$-open on $\pi'$. If $v_i=v_j$ then also $v_i$ is $C$-$\sigma'$-open on $\pi'$ (if $v_i$ is the left endnode or not). If $v_i \neq v_j$, then the new directed path $v_i \tuh \cdots \tuh v_j$ in $\pi'$ is $C$-$\sigma'$-open at $v_i$ because all nodes in between lie in the same stronly connected component $\Sc^G(v_i)$ (or $v_i$ is the left endnode anyways).

 Case (ii). Since case (i) is solved we can assume that we have $j<n$ with $v_j \hus v_{j+1}$ in $\pi$. 
 If $v_{i-1} \hut v_i$ on $\pi'$ (or $v_i$ the left endnode) then this case is analogous to case (i). 
 So we can also assume that we have $i > 0$ and $v_{i-1} \suh v_i$ on $\pi$. So $\pi$ looks as follows:
 \[\pi:\qquad \cdots v_{i-1} \suh v_i \sus \cdots \sus v_j \hus v_{j+1} \cdots. \]
 So there must be a smallest number $k \in \lC i,\dots,j\rC$ such that a collider appears at $v_k$ on $\pi$:
 \[\pi:\qquad \cdots v_{i-1} \suh v_i \tuh \cdots \tuh v_k \hus \cdots \sus v_j \hus v_{j+1} \cdots. \]
 Since $\pi$ is $C$-$\sigma'$-open we have $v_k \in \Anc^G(C)$. Since $v_i \in \Anc^G(v_k)$ (otherwise $v_k$ would not be the first collider appearing after $v_i$) we thus have that also $v_i \in \Anc^G(C)$.
 So if we replace the subwalk $v_i \sus \cdots \sus v_j$ of $\pi$ by the shortest directed path $v_i \hut \cdots \hut v_j$ in $G$ we then get for $\pi'$ the following situation:
 \[\pi':\qquad \cdots v_{i-1} \suh v_i \hut \cdots \hut  v_j \hus v_{j+1} \cdots, \]
 which is then $C$-$\sigma'$-open at $v_i$ as $v_i \in \Anc^G(C)$. Note that this holds also when $v_i=v_j$.
 If $v_i \neq v_j$ then $v_j$ is also $C$-$\sigma'$-open on $\pi'$ as $v_j$ points left to a node in the same strongly connected component as $v_j$. 

 So in all cases $\pi'$ stays $C$-$\sigma'$-open.
\end{proof}
\end{Lem}

\end{comment}


\begin{comment}
\begin{Lem}\label{lem:replace_walk}
Let $G=(J,V,E,L)$ be a CDMG, $C \subseteq V \cup J$ and $\pi = \lp  v_0 \sus \cdots \sus v_n \rp$ be a $C$-$\sigma'$-open walk in $G$. 
  Suppose $v_i \in \Sc^G(v_j)$ for some $i,j \in \{0,\dots,n\}$ with $i < j$.
We can then replace the subwalk $v_i \sus \cdots \sus v_j$ of $\pi$ by 
  \begin{enumerate}[label=(\roman*)]
      \item a shortest directed path $v_i \tuh \cdots \tuh v_j$ in $G$ if $j = n$ or if $v_j \tuh v_{j+1}$ on $\pi$, or 
      \item a shortest directed path $v_i \hut \cdots \hut v_j$ in $G$ otherwise,
  \end{enumerate}
that is entirely within $\Sc^G(v_j)$ and such that the modified walk $\pi'$ is still $C$-$\sigma'$-open.
\end{Lem}
\begin{proof}
$\pi'$ cannot become $C$-$\sigma'$-blocked at one of the initial nodes $v_0, \dots, v_{i-1}$ or at one of the final nodes $v_{j+1}, \dots, v_n$ on $\pi'$, since these nodes occur in the same local configuration on $\pi$ and are not $C$-$\sigma'$-blocked on $\pi$ by assumption.
Furthermore, $\pi'$ cannot become $C$-$\sigma'$-blocked through one of the nodes strictly between $v_i$ and $v_j$ on $\pi'$ (if there are any), since these nodes are all non-endnode non-colliders that only point to nodes in the same strongly connected component $\Sc^G(v_j)$.

Because $\pi$ is $C$-$\sigma'$-open, $v_j \notin C$ if $j = n$ or if $v_j \tuh v_{j+1}$ \PF{If $v_{j+1} \in \Sc^G(v_j)$ then $v_{j+1} \in C$ not necessarily excluded !!!} on $\pi$. 
This holds in particular in case (i). 
Similarly, $v_i \notin C$ if $i = 0$ or $v_{i-1} \hut v_i$ \PF{same here} on $\pi$.

In case (i), $\pi'$ is not $C$-$\sigma'$-blocked by $v_j$ because $v_j$ is a non-collider on $\pi'$ but $v_j \notin C$. 
Also $v_i$ does not $C$-$\sigma'$-block $\pi'$. Indeed, assume $v_i \ne v_j$ (otherwise there is nothing to prove).
If $i = 0$, or if $i > 0$ and $v_{i-1} \hut v_i$ on $\pi'$, then the same holds for $\pi$ and hence $v_i \notin C$; $v_i$ is then a non-collider on $\pi'$, but $v_i \notin C$. 
If $i > 0$ and $v_{i-1} \suh v_i$ on $\pi'$ then $v_i$ is a non-endpoint non-collider on $\pi'$ that does not point to a node in another strongly connected component.

Now consider case (ii). 
If $i = 0$ or $v_{i-1} \hut v_i$ on $\pi'$ then this case is analogous to case (i). 
So assume $i > 0$ and $v_{i-1} \suh v_i$ on $\pi'$. 
%If $v_i$ is an endpoint of $\pi'$, then $v_i = v_j$ and $k = n$ and therefore $v_j \notin C$, and hence $v_i$ and $v_j$ do not $C$-$\sigma'$-block $\pi'$. 
$v_i$ must be a collider on $\pi'$ (whether $v_i = v_j$ or not). 
Then on the subwalk of $\pi$ between $v_i$ and $v_j$ there must be a directed walk from $v_i$ to a collider that is in $\Anc^G(C)$, which implies that $v_i \in \Anc^G(C)$. Hence, $v_i$ will not $C$-$\sigma'$-block $\pi'$. 
Also $v_j$ cannot $C$-$\sigma'$-block $\pi'$.
Indeed, assume $v_i \ne v_j$ (otherwise there is nothing to prove). 
Since $v_j \hus v_{j+1}$ on $\pi'$, $v_j$ is a non-endpoint non-collider on $\pi'$ that does not point to a node in another strongly connected component.
\end{proof}
\end{comment}

\begin{comment}
\begin{Prp}\label{prp:sigma_opens-2}
  Let $G=(J,V,E,L)$ be a CDMG. For $C \subseteq J \cup V$, and $i, j \in J \cup V$, the following are
  equivalent:
  \begin{enumerate}
    \item there exists a $C$-$\sigma$-open walk between $i$ and $j$ in $G$;
    \item there exists a $C$-$\sigma'$-open walk between $i$ and $j$ in $G$;
    \item there exists a $C$-$\sigma'$-open path between $i$ and $j$ in $G$.
  \end{enumerate}
\end{Prp}
\begin{proof}
$1 \implies 2$: Suppose there exists a $C$-$\sigma$-open walk between $i$ and $j$.
%Then all non-colliders are either not in $C$, or only point towards nodes in the same strongly connected component;
%all colliders are in $C$. 
All colliders are in $C$, and hence in $\Anc^G(C)$. So this same walk is also $C$-$\sigma'$-open. 

$2 \implies 1$: Suppose there exists a $C$-$\sigma'$-open walk between $i$ and $j$. 
%Then all non-colliders are either not in $C$, or only point towards nodes in the same strongly connected component; 
Then all colliders are in $\Anc^G(C)$. If a collider is not in $C$, we can extend the walk by replacing this collider 
$\suh v \hus$ by the walk 
$\suh v \tuh \dots \tuh c \hut \dots \hut v \hus$ consisting of a directed walk starting from $v$ towards the first
node $c$ in $C$ it encounters, and from there tracing back the same walk in opposite direction to $v$. All nodes on this
walk except for $c$ ($v$ included) are non-colliders not in $C$. Extending this walk in such a way for each
collider not in $C$, we obtain an extended walk that is $C$-$\sigma$-open.

$2 \implies 3$:
%It suffices to show that if there exists a $C$-$\sigma$-open walk between $i$ and $j$ in $G$, there exists a $C$-$\sigma'$-open path between $i$ and $j$ in $G$. 
Let $\pi = (v_0,\dots,v_n)$ be a $C$-$\sigma'$-open walk between nodes $v_0$ and $v_0$ in $G$. 
If a node $w$ occurs more than once on $\pi$, let $v_i$ be the first node on $\pi$ and $v_j$ be the last node on $\pi$ that are in $\Sc^G(w)$.
We now use Lemma~\ref{lem:replace_walk} to construct a new walk $\pi'$ from $\pi$ by replacing the subwalk between $v_i$ and $v_j$ of $\pi$ by a particular directed path in $\Sc^G(w)$ between $v_i$ and $v_j$ in such a way that $\pi'$ is still $C$-$\sigma'$-open.
In $\pi'$, the number of nodes that occurs more than once is at least one less than in $\pi$, and all nodes within $\Sc^G(w)$ occur within a single segment.
This replacement procedure can be repeated until no nodes occur more than once. We have then obtained a
$C$-$\sigma'$-open path between $v_0$ and $v_n$.

$3 \implies 2$: trivial.
\end{proof}
Hence, it does not matter which of the three versions we use in the definition of $\sigma$-separation.
When checking $\sigma$-separation in a graph ``by hand'', we usually use the third formulation, since there are only a finite
number of paths to check. In proofs, though, it often is easier to make use of walks, since these can be concatenated into walks
(while one cannot in general concatenate two paths and again obtain a path).

\begin{Thm}
Let $G=(J,V,E,L)$ be a CDMG and $A,B,C \ins J \cup V$ subsets of nodes.
Then we have the equivalence:
\[A \Perp^\sigma_G B \given C  \qquad\iff\qquad A \Perp^{\sigma'}_G B \given C. \]
Similarly, we have the equivalence:
\[A \Perp^d_G B \given C  \qquad\iff\qquad A \Perp^{d'}_G B \given C. \]
\begin{proof}
    This follows from Proposition \ref{prp:sigma_opens} for $\sigma$-separation and analogously for d-separation. Details are left as an excercise.
\end{proof}
\end{Thm}

\end{comment}




\subsection{Acyclifications}\label{sec:graph_acyclifications}

It is possible to reformulate the notion of $\sigma$-separation in terms of d-separation on a modified and acyclic graph
by making use of the following construction, which will be the main tool to extend the acyclic theory to the cyclic one.
The construction was first proposed in the context of CBNs by \cite{Spirtes94,Spirtes95}.
\begin{Def}\label{def:cdg_acyclification}
    Given a CDMG $G = ( J, V, E, L)$, we call a CADMG $G' = (J', V', E', L')$
  an \emph{acyclification of $G$} if
  \begin{enumerate}[label=(\roman*)]
    \item $G'$ is acyclic;
    \item $G'$ has the same input nodes and output nodes as $G$, i.e.\ $J'=J$ and $V'=V$;
    \item for every pair of nodes $(i,j)$ such that $i \not\in\Sc^G(j)$:
      \begin{enumerate}
      \item $i \tuh j \in E'$ iff there exists a node $j' \in \Sc^G(j)$ such that $i \tuh j' \in E$;
      \item $i \huh j \in L'$ iff there exist nodes $i' \in \Sc^G(i), j' \in \Sc^G(j)$ such that $i' \huh j' \in L$;
      \end{enumerate}\label{def:acyclification_inter_scc}
  \item for every pair of distinct nodes $(i,j)$ such that $i \in\Sc^G(j)$: $i \tuh j \in E'$ or $i \hut j \in E'$ or $i \huh j \in L'$.\label{def:acyclification_intra_scc}
  \end{enumerate}
\end{Def}

The important property of acyclifications is that they can be used to express $\sigma$-separation in a (possibly cyclic) graph
in terms of $d$-separation in an acyclification.

\begin{Prp}\label{prp:sigma_sep_as_d_sep}
  Let $G = ( J, V, E, L)$ be a CDMG and $G'$ an acyclification of $G$. 
  Then for $A, B, C \subseteq V \cup J$ (not necessarily disjoint) subsets of nodes we have the equivalence:
  \[A \Perp^\sigma_G B \given C \iff A \Perp^\sigma_{G'} B \given C \iff A \Perp^d_{G'} B \given C.\]
\end{Prp}

\begin{proof}
We will show that there is a $C$-$\sigma$-open walk between $A$ and $B \cup J$ in $G$ if and only
if there is a $C$-$\sigma$-open walk between $A$ and $B \cup J$ in $G'$. Since $G'$ is acyclic, this
is in turn equivalent to the existence of a $C$-$d$-open walk between $A$ and $B \cup J$ in $G'$.

$\implies$: Suppose there is a $C$-$\sigma$-open walk $\pi = (v_0, \dots, v_n)$ between $A$ and $B \cup J$ in $G$. 
All its colliders are in $C$
and all its non-colliders are either not in $C$, or otherwise, point only to nodes in the same strongly connected component.
Note that each edge between two nodes in different strongly connected components in $G$ is also present in $G'$.
Edges between two nodes in the same strongly connected component, however, may not be present in $G'$.
Therefore, we will replace these edges with walks in $G'$. 
Consider a subwalk $(v_i,\dots,v_j)$ of maximum length that is entirely contained within a strongly connected component in $G$
(with possibly $i=j$).
We distinguish different cases and show for each case how this subwalk can be replaced by a subwalk in $G'$.
\begin{enumerate}[label=(\roman*)]
  \item 
$\suh v_i \cdots v_j \hus$: the subwalk between $v_i$ and $v_j$ has to contain a collider, say $w$,
which must be in $C$ since the walk between $v_i$ and $v_j$ is $C$-$\sigma$-open.
We can replace this subwalk by $\suh w \hus$ in $G'$ such that $w$ becomes a collider in $C$.

\item
$(\hut) v_i \cdots v_j \hus$:\footnote{We put parentheses around the first directed edge to indicate that
this case also applies if $v_i$ is an endnode, i.e., if $i = 0$.} 
here $v_i$ is a non-collider pointing to another strongly connected component or $v_i$ is an endnode,
and in both cases, $v_i \notin C$. Therefore, we can replace the subwalk by $(\hut) v_i \hus$ in $G'$,
such that $v_i$ becomes a non-collider not in $C$.

\item
$\suh v_i \cdots v_j (\tuh)$: analogous to the previous case, we can replace it by $\suh v_j (\tuh)$ in $G'$,
such that $v_j$ becomes a non-collider not in $C$.

\item
$(\hut) v_i \cdots v_j (\tuh)$: 
$v_i,v_j$ are both not in $C$ by assumption. If $i = j$, we replace this subwalk by $(\hut) v_i (\tuh)$ 
such that $v_i$ becomes a non-collider not in $C$. If $i < j$, we replace this subwalk by
$(\hut) v_i \sus v_j (\tuh)$ with $v_i \sus v_j$ any edge connecting $v_i$ and $v_j$ in $G'$,
such that both $v_i$ and $v_j$ become non-colliders not in $C$.
\end{enumerate}
By replacing all maximal subwalks of the original walk $\pi$ that are contained within a single strongly connected
component of $G$ in this way, we obtain a walk in the acyclification $G'$ that is $C$-$\sigma$-open by construction.
Note that the modified walk has the same endpoints ($v_0$ and $v_n$) as the original walk.

$\impliedby$: Suppose there is a $C$-$\sigma$-open walk $\pi'$ between $A$ and $B \cup J$ in $G'$.
All its colliders are in $C$, and all its non-colliders are not in $C$. 
We will construct a walk $\pi$ in $G$ with the same endpoints as $\pi'$ that is $C$-$\sigma$-open.

Consider a non-trivial subwalk $(v_i,\dots,v_j)$ on $\pi'$ of maximum length that is entirely contained within a strongly connected component of $G$. 
This subwalk may not be present in $G'$.
We distinguish different cases and show for each case how this subwalk can be replaced by a subwalk in $G$.
\begin{enumerate}[label=(\roman*)]
  \item 
$\suh v_i \cdots v_j \hus$: the subwalk between $v_i$ and $v_j$ has to contain a collider, say $w$,
which must be in $C$ since the walk between $v_i$ and $v_j$ is $C$-$\sigma$-open, and must be in
$\Sc^G(v_i) = \Sc^G(v_j)$ by assumption. We can replace this subwalk by
$\suh v_i \tuh \dots \tuh w \hut \dots \hut v_j \hus$ in $G$,
with possibly $v_i = w$ and possibly $w = v_j$, 
with all nodes in $\Sc^G(v_i)$. 
Note that the modified walk remains $C$-$\sigma$-open.

\item
$(\hut) v_i \cdots v_j \hus$:
here $v_i$ is a non-collider pointing to another strongly connected component or $v_i$ is an endnode,
and in both cases, $v_i \notin C$. 
We can replace this subwalk by a shortest directed walk
$(\hut) v_i \hut \dots \hut v_j \hus$ in $G$ with all nodes in $\Sc^G(v_i)$.
Note that the modified walk remains $C$-$\sigma$-open.

\item
$\suh v_i \cdots v_j (\tuh)$: analogous to the previous case, we can replace it by $\suh v_i \tuh \dots \tuh v_j (\tuh)$ in $G$.

\item
$(\hut) v_i \cdots v_j (\tuh)$: 
$v_i,v_j$ are both not in $C$ by assumption. We can replace this subwalk by a shortest directed walk
$(\hut) v_i \tuh \dots \tuh v_j (\tuh)$ in $G$ with all nodes in $\Sc^G(v_i)$.
The modified walk remains $C$-$\sigma$-open.
\end{enumerate}
In each of the four cases, in the modified walk both $v_i$ and $v_j$ become
either colliders in $C$, or non-colliders not in $C$, or non-colliders in $C$
that only point to a node in the same strongly connected component of $G$.

Now, we will replace edges on $\pi'$ between two strongly connected components that are not present in $G$.
For any directed edge $i \tuh j$ on $\pi'$ with $j \notin \Sc^G(i)$, there
must be a $j' \in \Sc^G(j)$ such that $i \tuh j'$ is present in $G$, and hence there must be a directed path
$j' \tuh \dots \tuh j$ entirely in $\Sc^G(j)$ such that we can use
$i \tuh j' \tuh \dots \tuh j$ as replacement in $G$ of the edge $i \tuh j$.
Similarly, for any bidirected edge $i \huh j$ on $\pi'$ with $j \notin \Sc^G(i)$, there
must be $i' \in \Sc^G(i)$ and $j' \in \Sc^G(j)$ such that $i' \huh j'$ is present in $G$, 
and hence there must be a walk $i \hut \dots \hut i' \huh j' \tuh \dots \tuh j$ in $G$, where
$i \hut \dots \hut i'$ is entirely in $\Sc^G(i)$ and $j' \tuh \dots \tuh j$ is entirely
in $\Sc^G(j)$, that we can use as replacement in $G$ of the edge $i \huh j$.
The new nodes introduced on $\pi$ in these replacements are non-colliders
that only point to nodes in the same strongly connected component. 
The endpoints of the replacement paths do not change their status:
if they were colliders in $C$ on $\pi'$ they still are 
on $\pi$, and if they were non-colliders not in $C$ on $\pi'$ they still are on $\pi$.

Hence we have constructed a walk $\pi$ in $G$ with the same endpoints as $\pi'$ that is $C$-$\sigma$-open.
\end{proof}

The following construction shows that acyclifications exist (but it is just one out of many possible ways to construct acyclifications).

\begin{Eg}[The standard acyclification]
    Let $G=(J,V,E,L)$ be a CDMG. Then we define the \emph{standard acyclification} of $G$ as the CDMG
    $ G'=(J,V,E',L')$ 
    where:
      \begin{align*} 
          E' := & \lC v_1 \tuh v_2  \st v_1 \in J \cup V, v_2 \in V, v_2 \notin \Sc^G(v_1),\exists v_2' \in \Sc^G(v_2):\, 
   v_1 \tuh v_2' \in E\rC,\\
          L' := &\lC v_1 \huh v_2  \st v_1,v_2 \in V, v_1 \neq v_2, \exists v_1' \in \Sc^G(v_1), v_2' \in \Sc^G(v_2):\, v_1 \huh v_2 \in L \rC \\ \cup &
\lC v_1 \huh v_2  \st v_1,v_2 \in V, v_1 \neq v_2, v_1 \in \Sc^G(v_2) \rC.
\end{align*}
The standard acyclification of a CDMG is acyclic, i.e.\ a CADMG.
\begin{proof}
    Assume that $G'$ is not acyclic. Then there exists a non-trivial cyclic directed walk in $G'$:
    \[ v_1 \tuh v_2 \tuh \cdots \tuh v_k \tuh v_1, \]
    for some $k \ge 1$. It is clear that $k \ge 2$ because clearly $v_1 \in \Sc^G(v_1)$, which rules out the existence of an edge $v_1 \tuh v_1 \in G'$. 
    For simplicity we now identify $v_{k+1} := v_1$ in the following.
    By construction of $G'$ for every $i=1,\dots,k$ there exists $v'_{i+1} \in \Sc^G(v_{i+1})$ 
    such that the edge $v_i \tuh v_{i+1}'$ exists in $G$. Since $v'_{i+1} \in \Sc^G(v_{i+1})$ there also exists a directed walk in $G$: 
     \[v'_{i+1} \tuh \cdots \tuh v_{i+1}.\]
     Concatenating all directed walks we get the cyclic directed walk in $G$:
    \[ v_1 \tuh v_2' \tuh \cdots \tuh v_2 \tuh v_3' \tuh \cdots \tuh v_k 
    \tuh v_1' \tuh \cdots \tuh v_1. \]
    This shows that $v_2 \in \Sc^G(v_1)$, which is a contradiction to the existence of the edge $v_1 \tuh v_2 \in G'$.
    So a non-trivial cyclic directed walk in $G'$ cannot exist in the first place. So $G'$ must be acyclic.
\end{proof}
\end{Eg}

\begin{figure}\label{fig:acyclifications}
  \centering
  \begin{tikzpicture}
    \begin{scope}
      \node[ndint] (X1) at (0,0) {$X_1$};
      \node[ndout] (X2) at (1.5,0) {$X_2$};
      \node[ndout] (X3) at (3,0) {$X_3$};
      \node[ndout] (X4) at (0,-1.5) {$X_4$};
      \node[ndout] (X5) at (3,-1.5) {$X_5$};
      \node[ndout] (X6) at (1.5,-2.5) {$X_6$};
      \node[ndint] (X7) at (-1.5,-1.5) {$X_7$};
      \node[ndout] (X8) at (0,-3) {$X_8$};
      \draw[arint] (X1) -- (X4);
      \draw[arout] (X2) -- (X4);
      \draw[arout] (X2) -- (X3);
      \draw[arout] (X3) -- (X5);
      \draw[arout] (X4) -- (X5);
      \draw[arout] (X5) -- (X6);
      \draw[arout] (X6) -- (X4);
      \draw[arint] (X4) -- (X8);
      \draw[arout] (X7) -- (X8);
    \end{scope}
    \begin{scope}[yshift=-5cm,xshift=-3.5cm]
      \node[ndint] (X1) at (0,0) {$X_1$};
      \node[ndout] (X2) at (1.5,0) {$X_2$};
      \node[ndout] (X3) at (3,0) {$X_3$};
      \node[ndout] (X4) at (0,-1.5) {$X_4$};
      \node[ndout] (X5) at (3,-1.5) {$X_5$};
      \node[ndout] (X6) at (1.5,-2.5) {$X_6$};
      \node[ndint] (X7) at (-1.5,-1.5) {$X_7$};
      \node[ndout] (X8) at (0,-3) {$X_8$};
      \draw[arint] (X1) -- (X4);
      \draw[arout] (X2) -- (X4);
      \draw[arout] (X2) -- (X3);
      \draw[arout] (X3) -- (X5);
      \draw[arout] (X4) -- (X5);
      \draw[arout] (X5) -- (X6);
      \draw[arout] (X4) -- (X6);
      \draw[arint] (X4) -- (X8);
      \draw[arout] (X7) -- (X8);
      \draw[arint] (X1) -- (X5);
      \draw[arint] (X1) -- (X6);
      \draw[arout] (X2) -- (X5);
      \draw[arout] (X2) -- (X6);
      \draw[arout] (X3) -- (X4);
      \draw[arout] (X3) -- (X6);
    \end{scope}
    \begin{scope}[yshift=-5cm,xshift=3.5cm]
      \node[ndint] (X1) at (0,0) {$X_1$};
      \node[ndout] (X2) at (1.5,0) {$X_2$};
      \node[ndout] (X3) at (3,0) {$X_3$};
      \node[ndout] (X4) at (0,-1.5) {$X_4$};
      \node[ndout] (X5) at (3,-1.5) {$X_5$};
      \node[ndout] (X6) at (1.5,-2.5) {$X_6$};
      \node[ndint] (X7) at (-1.5,-1.5) {$X_7$};
      \node[ndout] (X8) at (0,-3) {$X_8$};
      \draw[arint] (X1) -- (X4);
      \draw[arout] (X2) -- (X4);
      \draw[arout] (X2) -- (X3);
      \draw[arout] (X3) -- (X5);
      \draw[arlat] (X4) -- (X5);
      \draw[arlat] (X5) -- (X6);
      \draw[arlat] (X4) -- (X6);
      \draw[arint] (X4) -- (X8);
      \draw[arout] (X7) -- (X8);
      \draw[arint] (X1) -- (X5);
      \draw[arint] (X1) -- (X6);
      \draw[arout] (X2) -- (X5);
      \draw[arout] (X2) -- (X6);
      \draw[arout] (X3) -- (X4);
      \draw[arout] (X3) -- (X6);
    \end{scope}
  \end{tikzpicture}
  \caption{Top: CDMG $G$. Bottom: two acyclifications of $G$.}
\end{figure}


\subsection{Separoid Axioms for \textsigma-/d-Separation}

\begin{DefThm}[(Asymmetric) separoid axioms for $\sigma$-separation/d-separation]
    \label{d-sep-separoid-axioms}
    Let $G=(J,V,E,L)$ be a CDMG and $A,B,C,D \ins J \cup V$ subsets of nodes.
Then the ternary relations $\Perp = \Perp_G^d$ and $\Perp = \Perp_G^\sigma$ satisfy the following rules:
\begin{enumerate}[label=\alph*)]
    \item Extended Left Redundancy: 
    \item[] $D \ins A \implies D \Perp  B \given A$.
    \item $J$-Restricted Right Redundancy: 
    \item[] $A \Perp  \emptyset \given C \cup J$ always holds.
    \item $J$-Inverted Right Decomposition: 
    \item[] $A \Perp  B \given C \implies A \Perp  J \cup B \given C$.
	\item Left Decomposition:  
	\item[]$A \cup D \Perp  B \given C \implies D \Perp  B \given C$.
  \item Right Decomposition:  	
    \item[] $A \Perp  B \cup D\given C \implies A \Perp  D \given C$.
    \item Left Weak Union:
    \item[] $A \cup D \Perp  B  \given C \implies A \Perp  B \given D \cup C$.
  \item Right Weak Union:\item[] $A \Perp  B \cup D \given C \implies A \Perp  B \given D \cup C$.
\item Left Contraction:\item[] 
 $ (A \Perp  B \given D \cup C) \land (D \Perp  B \given C) \implies A \cup D \Perp  B \given C$.
\item Right Contraction:
\item[] $ (A \Perp  B \given D \cup C) \land (A \Perp  D \given C)  \implies A \Perp  B \cup D \given C$.
\item Right Cross Contraction:
\item[] $ (A \Perp  B \given D \cup C) \land (D \Perp  A \given C) \implies A  \Perp  B \cup D \given C$.
\item Flipped Left Cross Contraction: 
\item[] $ (A \Perp  B \given D \cup C) \land (B \Perp  D \given C) \implies B \Perp  A \cup D \given C$.
\end{enumerate}
In particular, we have the equivalences:
$$ (A \Perp  B \cup D \given C) \quad\iff\quad (A \Perp  B \given D \cup C) \quad\land\quad (A \Perp  D \given C)  ,$$
%and:
$$ (A \cup D \Perp  B \given C) \quad\iff\quad (A \Perp  B \given D \cup C) \quad\land\quad (D \Perp  B \given C).$$
%Finally, if $J = \emptyset$ then 
We also get:
\begin{enumerate}[resume,label=\alph*)]
    \item $J$-Restricted Symmetry: 
	\item[] $A \Perp  B \given C \cup J \implies B \Perp A \given C \cup J$. 
\end{enumerate}
For the special case of $J=\emptyset$ we have thus (unrestricted) Symmetry.
\end{DefThm}


\begin{Rem}
Let the assumptions be like in Theorem \ref{d-sep-separoid-axioms}.
    We also have the following rules: 
    \begin{enumerate}[resume,label=\alph*)]
        \setcounter{enumi}{12}
        \item Left Composition:
        \item[] $(A \Perp  B \given  C) \land (D \Perp  B \given C) \implies A \cup D \Perp  B \given C.$
        \item Right Composition:
        \item[] $(A \Perp  B \given  C) \land (A \Perp  D \given C) \implies A \Perp  B \cup  D \given C.$
         \item Left Intersection: If $A \cap D = \emptyset$ then:
        \item[] $(A \Perp  B \given D \cup  C) \land (D \Perp  B \given A \cup C) \implies A \cup D \Perp  B \given C.$       
        \item Right Intersection: If $B \cap D = \emptyset$ then:
        \item[] $(A \Perp  B \given D \cup  C) \land (A \Perp  D \given B \cup C) \implies A \Perp  B \cup D \given C.$
    \end{enumerate}
\end{Rem}
